\documentclass[11pt]{article}
\usepackage[T1]{fontenc}
\usepackage[utf8]{inputenc}
\usepackage{lmodern}
\usepackage[margin=1in]{geometry}
\usepackage{amsmath,amssymb,amsthm,mathtools,bm}
\usepackage{booktabs,array,longtable}
\usepackage{enumitem}
\usepackage{hyperref}
\newcommand{\OPHPaperReleaseID}{r1412}
\newcommand{\OPHPaperReleaseDate}{May 4, 2026}
\newcommand{\OPHPaperReleaseBanner}{%
\begin{center}
\small\textbf{Paper release:} \texttt{\OPHPaperReleaseID}\quad\textbf{Released:} \OPHPaperReleaseDate
\end{center}
}


\hypersetup{colorlinks=true,linkcolor=blue,citecolor=blue,urlcolor=blue}
\setlength{\emergencystretch}{3em}
\setlength{\tabcolsep}{5pt}
\renewcommand{\arraystretch}{1.08}

\newtheorem{theorem}{Theorem}[section]
\newtheorem{corollary}[theorem]{Corollary}
\newtheorem{proposition}[theorem]{Proposition}
\newtheorem{lemma}[theorem]{Lemma}
\newtheorem{definition}[theorem]{Definition}
\newtheorem{axiom}[theorem]{Axiom}
\newtheorem{assumption}[theorem]{Assumption}
\newtheorem{remark}[theorem]{Remark}

\newcommand{\SU}{\mathrm{SU}}
\newcommand{\U}{\mathrm{U}}
\newcommand{\Conf}{\mathrm{Conf}}

\title{Deriving the Particle Zoo from\\
Observer Consistency}
\author{B.\ M\"uller}
\date{\OPHPaperReleaseDate}

\begin{document}
\maketitle
\OPHPaperReleaseBanner

\begin{abstract}
Observer-Patch Holography (OPH) derives physics from consistency
conditions among overlapping observer patches on a regulated screen. This paper studies its
particle-spectrum branch. Starting from five axioms, the structural theorem package summarized in
\emph{Recovering Relativity and Standard Model Structure from Observer Overlap Consistency}~\cite{ophcompact},
and the dimensionless pixel constant
\[
P \equiv a_{\mathrm{cell}}/\ell_P^2,
\]
this paper asks how far one can derive the observed spectrum from one shared calibration
constant rather than from the large fitted parameter list of the Standard Model, typically counted
as \(19\) free parameters in the minimal massless-neutrino theory and at least \(26\) once
neutrino masses and mixing are included. On the derivation surface reported here, the photon, gluons,
and graviton are structural massless carriers; the electroweak calibration stage is closed on a
target-free source-only theorem that emits the \(W\) boson and \(Z\) boson; the Higgs/top critical
stage yields the Higgs boson and top quark;
and the quark branch emits continuation-level rows for the up quark, down quark, strange quark,
charm quark, and bottom quark. Charged leptons and flavor-labeled neutrinos remain
open for different reasons: the charged lane fixes the centered hierarchy but not the absolute
normalization, while the current weighted-cycle neutrino branch lands on the observed PMNS/hierarchy
pattern but still leaves one overall positive normalization scalar open. Hadrons remain execution-bound:
unlike the upstream elementary-particle lanes, the hadron branch requires large-scale classical
lattice-QCD computation and is not presently a realistic target for ordinary local hardware or
current cloud quantum devices. The purpose of the paper is therefore not to claim full closure, but to give a concise
end-to-end derivation map from observer-consistency axioms and \(P\) to the reported particle
outputs while marking every remaining closure object explicitly.
\end{abstract}

\section{Introduction}

Observer-Patch Holography (OPH) derives physics from observer consistency. Its
starting point is that physical data are organized in local patch algebras on a screen, nearby
observers must agree on overlaps, realized states obey a local maximum-entropy/refinement
principle, generalized entropy is recoverable, and the realized low-energy branch is selected by a
minimal admissibility criterion. From that basis one can derive, conditionally but explicitly, the
Lorentzian and Einstein branches together with the realized Standard Model gauge structure
\[
\frac{\SU(3)\times\SU(2)\times\U(1)}{\mathbb Z_6},
\qquad
N_g=3,
\qquad
N_c=3.
\]
The Lorentzian branch, Einstein recovery, and the realized Standard Model gauge structure written
above are derived in \cite{ophcompact}. That same paper also derives the hypercharge lattice and
the counting results \(N_g=3\), \(N_c=3\). The present paper uses those results as input and summarizes only the
downstream particle-spectrum continuation. If that derivation is correct, it should not stop at gauge structure. It should also constrain the
particle spectrum.

The particle branch studied here is a reverse-engineering problem. Instead of treating masses,
mixings, and Yukawa parameters as independent inputs, it asks how far one common dimensionless
pixel constant \(P\) can be propagated through the whole spectrum. The working hypothesis is that
one first infers \(P\) from a small electroweak calibration family and then uses the same \(P\)
through the bosonic, quark, charged-lepton, neutrino, and hadronic branches. The structural
massless carriers, the closed \(W\) boson/\(Z\) boson electroweak pair, the Higgs boson/top quark pair,
and several quark rows are emitted, while charged leptons, flavor-labeled neutrinos, and
hadrons remain non-closed, with hadrons specifically execution-bound. This paper develops that derivation chain from the axioms and \(P\) to the
reported particle outputs, while keeping the open closure boundary explicit.

\subsection{Why this particle zoo, and not another one?}

The shortest informal answer is that, in OPH, particles are not inserted by hand as a list of
elementary ingredients. They are the stable observer-visible excitations and carrier modes that
remain once three layers of structure have been fixed:
\begin{enumerate}[leftmargin=2em]
\item overlap consistency on the observer-patch network;
\item compact gauge reconstruction from refinement-stable edge-sector transport;
\item minimal admissible realization of the low-energy branch.
\end{enumerate}
The first layer says what kinds of local data can be compared consistently across neighboring
patches. The second says that persistent transportable sector data organize into a compact gauge
structure. The third selects which of the admissible low-energy branches is physically realized.
At that point one is no longer choosing an arbitrary particle catalog. One is reading off the
carrier and matter content of the realized branch.

That is why \emph{Recovering Relativity and Standard Model Structure from Observer Overlap Consistency}~\cite{ophcompact} is so important for this draft. That paper does
the structural work that fixes the realized low-energy package
\[
\frac{\SU(3)\times\SU(2)\times\U(1)}{\mathbb Z_6},
\qquad
N_g=3,
\qquad
N_c=3,
\]
together with one admissible Higgs doublet and the realized chiral matter content. In plain
language: OPH does not begin by guessing photons, gluons, quarks, leptons, and weak bosons and
then asking how to fit them. It first derives the realized gauge-and-matter branch, and that
branch determines which elementary carriers and matter families are available at low energy.

This also explains the phrase ``and no others'' in the only honest sense available at present.
The structural OPH claim is not that every imaginable beyond-Standard-Model extension has been
ruled out in every possible future continuation. The structural claim is narrower and stronger:
once one imposes the stated admissibility conditions, the realized low-energy branch is the
Standard-Model branch rather than a larger light-sector package. Additional light gauge bosons,
additional light Higgs multiplets, extra light chiral families, or low-energy supersymmetric
partners do not appear on that realized branch because they would enlarge the admissible package
that MAR selects against. In that sense the observed light particle zoo is not an arbitrary menu;
it is the minimal realized one.

The individual families then emerge for different reasons. The photon, gluons, and graviton are
structural carriers: once the realized electromagnetic, color, and dynamical-metric branches are
present, those massless carriers follow as symmetry-protected zeros. The weak bosons arise when
the electroweak gauge sector is calibrated quantitatively. Quarks and leptons arise from the
realized chiral matter package together with the three-generation structure \(N_g=3\). Their
family splittings are then read from the deeper overlap-transport and excitation machinery, not
from a second arbitrary postulate that says ``copy the family three times and assign masses by
hand.'' Hadrons come one layer later still: once color is realized and confined, stable hadrons
are the composite readout channels of the quark/gluon sector rather than separate fundamental
inputs.

The role of the pixel constant \(P\) is also important to state clearly. In OPH, \(P\) does not
decide whether the photon exists, whether there are three colors, or whether quarks come in three
generations. Those structural facts belong to the gauge-and-matter branch fixed before
calibration. What \(P\) does is set the shared quantitative scale on which the realized particle
content is read out numerically. So the logic of the paper is:
\[
\begin{aligned}
\text{observer consistency}
&\to
\text{realized gauge/matter branch}
\to
\text{particle content} \\
&\to
\text{common scale }P
\to
\text{quantitative spectrum}.
\end{aligned}
\]
That is the sense in which OPH aims to explain why a universe with the realized branch and the
shared pixel scale exhibits this particle zoo rather than some unrelated alternative one.

\section{Conceptual Overview: From \texorpdfstring{$S^2$}{S2} Overlap Obstructions to Particle Families}

Before the hard derivation chapters, it helps to say in plain but still technical language what
the OPH particle derivation is trying to do. The basic picture is not ``put particles on a
sphere.'' The picture is: start with a network of observer patches on a screen \(S^2\), ask how
their local descriptions can fail to glue perfectly, and then study which persistent obstruction
patterns survive repair, transport, and refinement. In OPH, those persistent patterns
are what eventually harden into gauge sectors, flavor structure, and particle families.

\subsection{Step 1: local observer patches do not automatically agree}

Imagine a finite network of observer patches covering the screen. Each patch carries local data,
and every neighboring pair of patches shares an overlap through which they can compare notes. If
all local descriptions matched perfectly on every overlap, then there would be no nontrivial
structure to study: the global state would be one big consistent codeword. But generic local
descriptions do \emph{not} line up so simply. They can disagree on overlaps, they can require
local repair maps, and they can carry nontrivial loop data when one transports information around
cycles in the patch network.

This is where the consensus paper's language of inconsistency potential, local repair, cycle
holonomy, and frustration becomes useful. Local agreement on one edge does not guarantee global
agreement around a loop. So the first genuinely interesting objects in the OPH story are not yet
particles. They are overlap defects, repair obstructions, and transport inconsistencies on the
\(S^2\) observer network.

\subsection{Step 2: stable overlap obstructions become sector data}

Once those overlap obstructions are organized algebraically, they stop looking like raw
mismatches and start looking like sector data. The collar and edge-center structure tells us that
boundary-visible information decomposes into sector labels, and the gauge-as-gluing perspective
reinterprets part of what looked like mismatch bookkeeping as redundancy or implementation
hiding. In other words, some of the nontrivial overlap structure is not ``noise'' but the
mathematical content that later becomes gauge structure.

At this stage the key idea is: if sector labels can be transported coherently under refinement,
and if their fusion and duality data satisfy the required categorical properties, then one can
reconstruct a compact gauge group from them. This is the first big jump from screen-network
obstructions to familiar particle-physics structure. The obstruction story has been compressed
into transportable edge sectors, and those edge sectors become the raw material for gauge
symmetry.

\subsection{Step 3: the realized Standard Model branch fixes the structural carrier skeleton}

Once compact gauge reconstruction is in place, MAR selects the realized Standard Model branch.
That is where the particle-spectrum story begins to resemble ordinary particle physics. The
realized branch fixes
\[
\frac{\SU(3)\times\SU(2)\times\U(1)}{\mathbb Z_6},
\qquad
N_g=3,
\qquad
N_c=3,
\]
and from there several carrier statements are structural. The photon, gluons, and
graviton are not late numerical accidents. They are the symmetry-protected massless carriers of
the realized electromagnetic, color, and dynamical-metric branches. So one direct route from
screen-network obstruction language to particle content is:
\[
\begin{aligned}
\text{overlap obstruction}
&\to \text{edge-sector transport}
\to \text{compact gauge reconstruction} \\
&\to \text{realized Standard Model branch}
\to \text{structural carrier content}.
\end{aligned}
\]

\subsection{Step 4: refinement-stable transport on the three-generation bundle becomes flavor structure}

The next jump is subtler and more interesting. The Standard Model branch tells us that three
generations are realized, but it does not yet tell us what distinguishes one family from another.
That is where the flavor derivation enters. Its starting point is still obstruction
language, but in a more refined form: family transport kernels, generation-bundle branch
generators, same-label eigenline transport, overlap-edge cocycles, projector data, spectral
gaps, pair suppressions, and cycle phases.

Semi-informally, the idea is that persistent transport defects on the three-generation charged
bundle organize themselves into a family observable. Once that object exists, it can be pushed
forward into common sector-response data for the up/down quark sectors, the charged-lepton
sector, and the neutrino sector. In that sense, the particle families are not attached to the
screen by hand. They are supposed to emerge from how the refinement-stable observer network
stores and transports nontrivial overlap structure across the realized three-generation bundle.

\subsection{Step 5: the same family backbone branches into different particle families}

Once a common family object exists, the downstream sectors stop being arbitrary lists and become
different readout branches of the same deeper transport structure.

\paragraph{Bosons.}
The bosonic branches split early. Structural massless carriers follow from the realized
gauge and gravity branches. The \(W\) boson and \(Z\) boson then appear on the electroweak
calibration stage, while the Higgs boson appears on the Higgs/top critical stage.

\paragraph{Quarks.}
Quarks are the first matter family with reported numerical rows. The quark derivation
takes the shared flavor data, performs the mean split and sector descent, builds forward Yukawa
surfaces, and emits continuation-level reported rows for the up quark, down quark, strange quark,
charm quark, and bottom quark. The top quark is read out on the Higgs/top critical
stage instead.

\paragraph{Charged leptons.}
Charged leptons are no longer described here as a Koide-assisted fit. The exact statement is
narrower and better. On the ordered charged carrier
\[
(-1,x_2,1),
\qquad
x_2=-0.5175863354681689,
\]
every charged source pair \((\eta,\sigma)\) determines the centered charged logs by
\[
e_{\log,\mathrm{centered}}
=
-\frac{(3+x_2)\sigma-\eta}{6},
\qquad
\mu_{\log,\mathrm{centered}}
=
\frac{x_2\sigma-\eta}{3},
\qquad
\tau_{\log,\mathrm{centered}}
=
\frac{(3-x_2)\sigma+\eta}{6}.
\]
So the centered charged hierarchy is closed once a charged source pair exists, but the numerical
masses are not. The present theorem-grade waiting set is
\[
\widehat C_e^{\mathrm{cand}}
\xrightarrow{\ \texttt{oph\_generation\_bundle\_branch\_generator\_splitting}\ }
\widehat C_e
\quad\text{then}\quad
A_{\mathrm{ch}}
\quad\Longrightarrow\quad
(g_e,\Delta_e^{\mathrm{abs}}).
\]
Here \(\widehat C_e^{\mathrm{cand}}\) is the latent charged sector-response candidate, and the
live corpus proves neither exact vanishing nor uniform quadratic smallness of the descended
commutator after the central split. On the absolute side, the equalizer route is a no-go under
common-shift symmetry; the future slot is one affine-covariant anchor
\(A_{\mathrm{ch}}(\ell + c\,\mathbf 1)=A_{\mathrm{ch}}(\ell)+c\).
A minimal operator-side extension route is a centered Schur-type \(P\!\to\!Q\!\to\!P\) transfer
theorem for the descended commutator after the central split; that route would promote
\(\widehat C_e^{\mathrm{cand}}\), but it would still leave \(A_{\mathrm{ch}}\) as a separate
missing object.
The clean future absolute realization is an uncentered charged response lift carrying a
determinant line, in which
\[
A_{\mathrm{ch}}=\tfrac13 \log \det(Y_e)=\tfrac13 \operatorname{tr}(\log Y_e).
\]
The strongest continuation bridge is the balanced Hermitian circulant / \(Q=2/3\) / \(\delta=2/9\) branch, which gives near-exact centered-log closure at residual norm
\[
\left\|E_{\log,\mathrm{centered}}^{\mathrm{cont}}
-
E_{\log,\mathrm{centered}}^{\mathrm{ref}}\right\|
\approx
2.13\times10^{-5},
\]
but that bridge is compare-only and not an emitted mass theorem.

\paragraph{Neutrinos.}
The current weighted-cycle / Majorana-holonomy branch closes the PMNS/hierarchy shape and leaves no further discrete branch object on that lane. The remaining freedom is the positive rescaling orbit of the absolute spectrum, so one positive mass-normalization scalar remains open. Charged leptons and neutrinos are therefore open for different reasons: charged leptons still miss promotion of \(\widehat C_e^{\mathrm{cand}}\) plus one affine-covariant absolute anchor \(A_{\mathrm{ch}}\), while neutrinos miss one positive normalization scalar after the shape closure is already in hand.

\paragraph{Status boundary.}
The honest Stage-5 headline is therefore this: quarks have continuation-level rows; charged leptons have an exact same-carrier readback theorem together with an exact absolute-scale obstruction; neutrinos have a current weighted-cycle continuation with one scalar still open. None of those later-family statements should be promoted to the recovered-core surface fixed earlier by the gauge and relativity chain.
\subsection{Structural results and open closures}

The move from overlap structure to compact gauge reconstruction
and the realized Standard Model branch belongs to the OPH structural core. The move from
the realized branch to structural carriers and to the electroweak/Higgs-top bosonic outputs is
reflected in the reported numerical rows. The move from the flavor obstruction story to full
family closure is only partly complete: quarks have continuation-level rows,
charged leptons and neutrinos do not, and hadrons remain behind an explicit nonperturbative
frontier.

That is exactly why this section belongs near the front of the paper. It gives the reader the
interesting conceptual picture first: particle families are supposed to be the stabilized and
readable descendants of transport obstructions on an observer network over \(S^2\). The later
chapters then say, sector by sector, how far that picture has been turned into explicit
mathematics and reported outputs.

\section{Calibration and the Reverse-Engineering Strategy}

Before the status table, it is worth stating what the reverse-engineering claim actually is. In
ordinary phenomenological use, the Standard Model does not derive the observed particle masses and
mixings from one common microscopic constant. It is usually presented with \(19\) free parameters
in the minimal massless-neutrino theory, and with at least \(26\) once neutrino masses and mixing
are included, depending on neutrino-sector conventions. The OPH particle derivation is trying to
compress that situation radically. The aim is not to feed in one Yukawa coupling, one mass, or
one mixing angle per sector, but to infer one universal pixel constant \(P\) from a small
calibration family and then demand that the same \(P\) drive all downstream bosonic, quark,
lepton, neutrino, and hadronic branches.

That statement needs one immediate clarification. The present OPH formulation still carries the
screen-capacity input \(N_{\mathrm{scr}}\) in its cosmological branch. So the claim here is
not that every continuous number appearing anywhere in the full OPH framework has disappeared. The
claim is narrower and more interesting: the particle-spectrum reverse-engineering derivation is
organized around one common nontrivial scale input \(P\), rather than a long particle-by-particle
parameter list.

\subsection{Why calibration is necessary}

Semi-informally, the structural OPH theorems tell us \emph{what kind} of particle world is
realized, but not yet \emph{what absolute scale in GeV} that world occupies. The overlap,
modular, gauge, anomaly, and admissibility arguments recover the realized Standard Model branch,
the exact hypercharge pattern, three generations, three colors, and the symmetry-protected
massless carriers. But those structural results do not by themselves tell us the numerical values
of the \(W\) boson mass, the Higgs boson mass, or the up-quark mass. Calibration is the bridge
from the screen-side structural world to observed laboratory units.

Mathematically, the calibration problem can be phrased as follows. Let
\[
\mathcal C=\{O_i^{\mathrm{obs}}\}
\]
be a small family of observed calibration quantities, and let
\[
O_i^{\mathrm{pred}} = F_i(P)
\]
be the forward OPH prediction for each such quantity when the pixel constant is set to \(P\).
Then each calibration observable defines an implied pixel constant
\[
P_i := F_i^{-1}\!\bigl(O_i^{\mathrm{obs}}\bigr),
\]
whenever the inverse is locally meaningful. Exact single-\(P\) closure means that the same value
\[
P_* \quad \text{satisfies} \quad P_i=P_* \ \text{for all calibration observables}.
\]
Once such a common \(P_*\) exists, the real predictive question becomes whether all downstream
particle readouts can be written as
\[
X_j = G_j(P_*)
\]
with no new sector-specific fit constants inserted by hand.

\subsection{How electroweak calibration works}

In the present implementation, the calibration family is electroweak. The electroweak calibration
stage
tracks the observable set
\[
\{m_W,\ m_Z,\ \alpha_{\mathrm{em}}(m_Z),\ \sin^2\theta_W(m_Z),\ v\},
\]
and asks whether one and the same \(P\) can support them all. The construction proceeds as follows: start from \(P\), build the single-\(P\) running
electroweak family, reduce it to the two-scalar carrier \((\sigma_{\mathrm{EW}},
\eta_{\mathrm{EW}})\), apply the carrier selector on that reduced chart, and read off the
\(W\) boson/\(Z\) boson mass pair from the selected carrier point. The Higgs/top critical stage
then inherits this calibrated electroweak core rather than reopening a new
free-input sector.

This is why calibration is not a side detail. It is the place where the particle derivation shows
whether it is really doing reverse engineering or merely hiding many fitted inputs behind fancy
notation. In the electroweak branch, the public \(W/Z\) rows are emitted by a target-free
source-only repair theorem from the calibrated source basis. The derivation also contains two
auxiliary validation surfaces: the exact current-carrier chart and the freeze-once coherent repair
surface. On the calibration tier, the electroweak mass-side lane has no open blocker.

Throughout this paper, labels such as \(D10\), \(D12\), \(D12_{ud}^{\mathrm{mass}}\), and
\(\mathcal R_{D12}^{ud}\) are internal OPH ledger labels inherited from
\cite{ophcompact}, not external references. In that ledger, \(D10\) names the Phase-II
electroweak calibration branch and \(D12\) names the Phase-III phenomenological-continuation
branch. The current public paper ledger has no separate D11 row. The compact paper's
overview/dependency table is the canonical key, and this paper uses those labels only to locate
each particle-side object inside the broader OPH dependency map.

\subsection{What electroweak calibration implies}

Calibration fixes the common pixel constant \(P\) on the declared electroweak
running/matching surface and thereby fixes the source basis
\((\alpha_{2,m_Z},\alpha_{Y,m_Z},\eta_{\mathrm{source}},v)\). The target-free
source-only repair theorem emits one coherent quintet
\((W,Z,\alpha_{\mathrm{em}}^{-1},\sin^2\theta_W,v)\) from that calibrated basis. The
derivation also includes two explicit validation surfaces beneath the public theorem output: the
exact current-carrier chart and the freeze-once coherent repair surface. The same tiered logic is
useful later in the paper because flavor and neutrino derivations do have open theorem objects,
whereas the electroweak calibration lane does not.

\section{Results at a Glance}

In plain language, the particle derivation has three different kinds of rows: rows that are
structurally fixed, rows that are numerically emitted, and rows whose chains stop short of a promotable mass prediction. This section summarizes that split.

This section should be read by claim tier rather than by raw numerical proximity. The same absolute residual can mean either a closed emitted value or a compare-only continuation benchmark. The closed emitted-output set is presently the electroweak calibration pair \(m_W\) and \(m_Z\). Charged-lepton near-matches remain compare-only continuation diagnostics, and the current weighted-cycle neutrino branch remains one positive normalization scalar short of promoted absolute masses.

\begin{table}[h]
\centering
\small
\begin{tabular}{@{}>{\raggedright\arraybackslash}p{0.18\textwidth}>{\raggedright\arraybackslash}p{0.15\textwidth}>{\raggedright\arraybackslash}p{0.22\textwidth}>{\raggedright\arraybackslash}p{0.16\textwidth}>{\raggedright\arraybackslash}p{0.20\textwidth}@{}}
\toprule
Sector & Status & Derivation stage & Emitted value(s) & Remaining step \\
\midrule
Structural carriers & structural & gauge/gravity structural theorems & photon, gluons, graviton & none on the particle side \\
Electroweak bosons & calibration & electroweak calibration stage & \(W\) boson, \(Z\) boson & no open mass-side blocker on the active calibration surface \\
Higgs/top stage & secondary quantitative & Higgs/top critical stage & Higgs boson, top quark & presentation-level rigor and synchronization of the closed forward seed \\
Quark family & continuation & forward Yukawa derivation & up, down, strange, charm, bottom quarks & emitted same-family mass ray \(D12_{ud}^{\mathrm{mass}}\), one intrinsic scale-fixing law on that ray, branch promotion beyond the current D12 ray, Yukawa dictionary \\
Charged leptons & continuation & ordered-log and support-extension derivation & no promoted masses & centered charged shape is closed, but the theorem emits only the quotient class modulo common shift; theorem-grade waiting set is promotion of the latent candidate \(\widehat C_e^{\mathrm{cand}}\) via branch-generator splitting, then one affine-covariant absolute anchor \(A_{\mathrm{ch}}\) \\
Neutrinos & continuation & weighted-cycle PMNS/hierarchy closure & no promoted masses & one positive neutrino mass-normalization scalar remains open; absolute masses and absolute \(\Delta m^2\) are compare-only until it is emitted \\
Hadrons & execution-bound nonperturbative continuation & stable-channel correlator and readout derivation & no paper-derived hadron masses & backend correlator dump, executed unquenching, and production systematics \\
\bottomrule
\end{tabular}
\end{table}

Two points are worth stating explicitly. First, the only matter-family rows with emitted
numerical values are the quark rows, and even those remain continuation-level rather
than theorem-closed. Second, ``the bosons are finished'' is too blunt. The photon, gluons, and
graviton are structural zeros; the \(W\) boson and \(Z\) boson come from electroweak calibration; and
the Higgs boson sits on the Higgs/top critical stage. Those are all real outputs, but they live at
different tiers.

This is the status boundary that governs the rest of the paper. Whenever a later chapter becomes
detailed about an open family, it should still be read through this table: the point is to
document the constructive chain faithfully, not to silently promote it.

\paragraph{Strongest current numerical matches.}
For quick quantitative orientation, the strongest current rows are:

\begin{table}[h]
\centering
\scriptsize
\setlength{\tabcolsep}{3pt}
\begin{tabular}{@{}>{\raggedright\arraybackslash}p{0.20\textwidth}>{\raggedright\arraybackslash}p{0.14\textwidth}>{\raggedright\arraybackslash}p{0.17\textwidth}>{\raggedright\arraybackslash}p{0.12\textwidth}>{\raggedright\arraybackslash}p{0.08\textwidth}>{\raggedright\arraybackslash}p{0.22\textwidth}@{}}
\toprule
Particle / observable & OPH value & PDG / reference & Error & Unit & Status / caveat \\
\midrule
Photon & \(0\) & \(<10^{-27}\) & within bound & GeV & exact structural zero; compared against the current photon mass upper bound \\
Gluon & \(0\) & no direct free-particle mass measurement & \(n/a\) & GeV & exact structural zero for the color gauge sector; free gluons are confined \\
Graviton & \(0\) & \(<10^{-32}\) & within bound & GeV & exact structural zero; compared against the current graviton mass upper bound \\
\(W\) boson & \(80.37700001539531\) & \(80.377\) & \(+1.54\times 10^{-8}\) & GeV & emitted on the closed target-free electroweak calibration theorem \\
\(Z\) boson & \(91.18797807794321\) & \(91.1879781\) & \(-2.21\times 10^{-8}\) & GeV & emitted on the same closed electroweak calibration theorem \\
Higgs boson & \(125.218922\) & \(125.19953\) & \(+0.019392\) & GeV & secondary quantitative row on the Higgs/top critical stage \\
top quark & \(172.388646\) & \(172.352355\) & \(+0.036291\) & GeV & secondary quantitative row on the Higgs/top critical stage \\
Neutrino solar splitting \(\Delta m_{21}^2\) & \(7.489806641884242\times10^{-5}\) & \(7.49\times10^{-5}\) & \(-1.93\times10^{-9}\) & \(\mathrm{eV}^2\) & compare-only absolute splitting; \(\lambda_\nu\) still open \\
Neutrino atmospheric splitting \(\Delta m_{32}^2\) & \(2.438\times10^{-3}\) & \(2.438\times10^{-3}\) & \(0\) & \(\mathrm{eV}^2\) & compare-only atmospheric anchor used to place the branch on an eV scale \\
Neutrino normal-order splitting \(\Delta m_{31}^2\) & \(2.5128980664188426\times10^{-3}\) & \(2.5129\times10^{-3}\) & \(-1.93\times10^{-9}\) & \(\mathrm{eV}^2\) & compare-only absolute splitting; \(\lambda_\nu\) still open \\
Neutrino hierarchy ratio \(\Delta m_{21}^2/\Delta m_{32}^2\) & \(0.030721110097966534\) & \(0.030721903199343724\) & \(-7.93\times10^{-7}\) & dimensionless & theorem-grade scale-free hierarchy ratio on the weighted-cycle branch \\
\bottomrule
\end{tabular}
\end{table}

\subsection{Detailed particle inventory}

\begin{longtable}{@{}>{\raggedright\arraybackslash}p{0.14\textwidth}>{\raggedright\arraybackslash}p{0.15\textwidth}>{\raggedright\arraybackslash}p{0.14\textwidth}>{\raggedright\arraybackslash}p{0.22\textwidth}>{\raggedright\arraybackslash}p{0.27\textwidth}@{}}
\toprule
Family & Particle & Status & OPH value or strongest derived benchmark & Remaining step \\
\midrule
\endfirsthead
\toprule
Family & Particle & Status & OPH value or strongest derived benchmark & Remaining step \\
\midrule
\endhead
\bottomrule
\endfoot
Structural carriers & photon & structural & \(0\) & none \\
Structural carriers & gluons & structural & \(0\) for the color gauge sector & none \\
Structural carriers & graviton & structural & \(0\) for the dynamical-metric spin-2 carrier & none \\
Electroweak calibration & \(W\) boson & calibration & \(80.37700001539531~\mathrm{GeV}\) from the target-free source-only repair theorem on the calibrated D10 source basis & no open D10 mass-side blocker; machine-scale agreement with the freeze-once coherent repair validation surface \\
Electroweak calibration & \(Z\) boson & calibration & \(91.18797807794321~\mathrm{GeV}\) from the same closed theorem on the calibrated D10 source basis & same closed target-free electroweak theorem as the \(W\) row \\
Higgs/top critical stage & Higgs boson & secondary quantitative & \(125.218922~\mathrm{GeV}\) & present a proof package for the closed one-scalar forward seed \\
Quark family & top quark & secondary quantitative & \(172.388646~\mathrm{GeV}\) from the Higgs/top critical stage & same one-scalar seed proof package as the Higgs row \\
Quark family & up quark & overlap continuation & reported value \(0.00228766~\mathrm{GeV}\); stronger two-scalar candidate \(0.00215992~\mathrm{GeV}\) & the current D12 sheet is a strict no-go for the physical CKM shell; emit one discrete left-handed same-label relative-sheet selector \(\sigma_{ud}\), then derive the mass-side scale law on the selected branch \\
Quark family & down quark & overlap continuation & reported value \(0.00512144~\mathrm{GeV}\); stronger two-scalar candidate \(0.00470005~\mathrm{GeV}\) & same wrong-branch D12 frontier; same-sheet CKM/CP transport is closed but cannot be repaired by rephasing or scalar rescaling \\
Quark family & strange quark & overlap continuation & reported value \(0.0916111~\mathrm{GeV}\); stronger two-scalar candidate \(0.0935045~\mathrm{GeV}\) & same wrong-branch D12 frontier; the next honest quark object is the discrete left-handed same-label selector \(\sigma_{ud}\) rather than a continuous mass scalar \\
Quark family & charm quark & overlap continuation & reported value \(1.256692~\mathrm{GeV}\); stronger two-scalar candidate \(1.272939~\mathrm{GeV}\) & same wrong-branch D12 frontier; selected-branch mass-scale fixing is downstream of branch choice \\
Quark family & bottom quark & overlap continuation & reported value \(3.820118~\mathrm{GeV}\); stronger two-scalar candidate \(4.182754~\mathrm{GeV}\) & same wrong-branch D12 frontier; CKM/CP closure on the current sheet does not by itself identify the physical quark branch \\
Charged leptons & electron & continuation gap & \(n/a\); continuation pair \((\eta,\sigma)=(-6.72959,\,8.15406)\) gives near-exact centered-log shape closure, but the absolute scale remains quotient-only; compare-only target \(g_e^\star=0.0457789\), equivalently \(\Delta_e^{\mathrm{abs},\star}=3.00398633\) & promote the latent candidate \(\widehat C_e^{\mathrm{cand}}\) by closing branch-generator splitting, then derive a theorem-grade affine-covariant anchor \(A_{\mathrm{ch}}\), then read out \(g_e=e^{A_{\mathrm{ch}}}\) and \(\Delta_e^{\mathrm{abs}}=\log g_{\mathrm{ch,shared}}-A_{\mathrm{ch}}\) \\
Charged leptons & muon & continuation gap & \(n/a\); same centered-shape closure and common-shift no-go as the electron row & same \(\widehat C_e^{\mathrm{cand}} \rightarrow\) promotion \(\rightarrow A_{\mathrm{ch}} \rightarrow g_e\) closure chain \\
Charged leptons & tau lepton & continuation gap & \(n/a\); same centered-shape closure and common-shift no-go as the electron row & same \(\widehat C_e^{\mathrm{cand}} \rightarrow\) promotion \(\rightarrow A_{\mathrm{ch}} \rightarrow g_e\) closure chain \\
Neutrinos & \(\nu_e\) & weighted-cycle continuation & \(n/a\); current weighted-cycle branch emits \(\theta_{12}=34.2259^\circ\), \(\theta_{23}=49.7228^\circ\), \(\theta_{13}=8.68636^\circ\), \(\delta=305.581^\circ\), \(J=-0.02753\), and \(\Delta m_{21}^2/\Delta m_{32}^2=0.03072111\); with compare-only anchor \(\Delta m_{32}^2=2.438\times10^{-3}\,\mathrm{eV}^2\), it gives \(m_1=0.01746\), \(m_2=0.01948\), \(m_3=0.05308~\mathrm{eV}\) & emit the one positive neutrino mass-normalization scalar; no hidden discrete branch remains on that lane, so the absolute masses and absolute \(\Delta m^2\) stay anchor-based only because of the open scale \\
Neutrinos & \(\nu_\mu\) & weighted-cycle continuation & \(n/a\); same weighted-cycle branch as the electron-neutrino row & same one-scalar normalization frontier \\
Neutrinos & \(\nu_\tau\) & weighted-cycle continuation & \(n/a\); same weighted-cycle branch as the electron-neutrino row & same one-scalar normalization frontier \\
Hadrons & proton & execution-bound & not a paper-derived proton row; surrogate finest-ensemble \(N_{\mathrm{iso}}\) candidate \(0.938960210578~\mathrm{GeV}\) and production geometry summary with 3 ensembles / 6 cfg & replace the surrogate bridge with the backend correlator dump from real production execution and production systematics \\
Hadrons & neutron & execution-bound & not a paper-derived neutron row & production unquenched execution plus isospin-resolved hadron systematics \\
Hadrons & neutral pion proxy & execution-bound & not a paper-derived pion row; surrogate finest-ensemble \(\pi_{\mathrm{iso}}\) candidate \(0.135039383836~\mathrm{GeV}\) & replace the surrogate bridge with the backend correlator dump from real production execution and production systematics \\
Hadrons & \(\rho(770)^0\) proxy & execution-bound & not a paper-derived resonance row & production execution plus finite-volume resonance extraction \\
\end{longtable}

\section{Axioms, Screen Architecture, Inputs, and Theorem Package}

In plain terms, this chapter fixes the rules of the game before we start talking about particles.
The particle derivation uses the same five axioms, the same two declared external continuous
inputs, and the same structural theorem chain that yields relativity and Standard Model
structure in OPH. What matters here is not the publication history of those
results but their logical role in the particle argument: first fix the axioms and the screen-side
regulated picture, then state the external inputs, and then record exactly which theorem packages
are being used later in the derivation.

\subsection{Canonical OPH basis}

The particle-spectrum derivation uses the same five axioms as the Standard-Model/relativity derivation paper~\cite{ophcompact}. They are not
reintroduced here as a new package; they are repeated because every later particle-family
chapter depends on them either directly or through the structural OPH chain summarized in the
companion paper.

\begin{axiom}[Screen Net]\label{ax:screen}
Physical data is organized on a horizon screen \(S^2\) carrying a net of local algebras
\[
P \longmapsto \mathcal A(P)
\]
for connected patches \(P\subset S^2\), with isotony
\[
P\subset Q \implies \mathcal A(P)\subset \mathcal A(Q).
\]
\end{axiom}

\begin{axiom}[Overlap Consistency]\label{ax:overlap}
For overlapping patches \(P_1\cap P_2\neq\varnothing\), the local states induced on the shared
algebra agree:
\[
\omega_{P_1}|_{\mathcal A(P_1\cap P_2)}
=
\omega_{P_2}|_{\mathcal A(P_1\cap P_2)}.
\]
\end{axiom}

\begin{axiom}[Local MaxEnt and Refinement Stability]\label{ax:maxent}
At the regulator scale \(\ell_{\mathrm{UV}}\), the realized branch is selected by maximizing
entropy subject to a finite family of gauge-invariant local constraints
\[
\mathcal C_{\ell_{\mathrm{UV}}}=\{O_a(x)\}_{a=1}^{N_{\mathrm{con}}},
\]
with support radius \(O(\ell_{\mathrm{UV}})\) and with label set independent of the number of
regulator cells. Under refinement, the same finite constraint family persists, so the realized
states belong to one common finite-dimensional MaxEnt family rather than unrelated maximizers.
\end{axiom}

\begin{axiom}[Recoverable Generalized Entropy]\label{ax:entropy}
A generalized entropy functional exists on caps,
\[
S_{\mathrm{gen}}(C)=S_{\mathrm{bulk}}(C)+\langle L_C\rangle,
\]
where \(L_C\) is a positive edge-center entropy functional and the semiclassical branch
identifies its leading coarse-grained contribution with \(A(\partial C)/(4G)\). The functional
obeys the recoverability and focusing structure required by the collar and null-modular
arguments.
\end{axiom}

\begin{axiom}[Minimal Admissible Realization]\label{ax:mar}
Among admissible realized low-energy sector packages \(\mathfrak S\) consisting of the connected
Lie gauge-sector image relevant in the EFT branch, its admissible light chiral matter content,
and one Higgs doublet, the realized package is lexicographically minimal under
\[
C(\mathfrak S)=\bigl(\chi_{\mathrm{cpl}},\,N_{\mathrm{nonab}},\,N_c,\,N_g\bigr),
\]
subject to loop coherence, anomaly freedom, refinement-stable light chiral matter,
single-Higgs Yukawa completable structure with one connected abelian charge factor, intrinsic
CP capability, and weak-sector UV completeness.
\end{axiom}

For particle physics, the first four axioms provide the observer-centric kinematic and
entropic/modular background. The fifth axiom, MAR, is the selector that turns ``some compact
gauge group reconstructed from the refinement-stable edge-sector category'' into the realized
Standard Model branch. The particle paper therefore uses MAR constantly, but only after the
compact-gauge reconstruction step has been separated from the later realized-branch
selection.

\subsection{Regulated screen architecture and patch language}

The paper \emph{Recovering Relativity and Standard Model Structure from Observer Overlap Consistency}
states the axioms in algebraic language. The screen-microphysics work
paper supplies one concrete regulated realization of that language: a cellulated horizon screen,
finite link or quantum-link registers on the edges, record registers on selected vertices or
coarse cells, boundary-fixed patch algebras, overlap algebras on collars, typed readout packets,
and local repair instruments. This regulated architecture is important for the particle derivation
for three reasons.

First, it makes the screen picture operational. A patch is not a metaphor but a finite local
algebra carried by one region of the cellulated screen, and an overlap is not a slogan but a
boundary-visible algebra with declared shared observables. Second, it makes the edge-sector
language concrete. The particle derivation ultimately uses edge sectors, transport, refinement,
fusion, and overlap data to build the gauge and flavor branches; the microphysics paper shows
how those objects can be realized at fixed cutoff by finite gauge-aware registers.
Third, it clarifies the measurement interface. The observer-facing measurement package is not an
added Copenhagen layer but a theorem-bearing fixed-cutoff statement about a central record
algebra, Born probabilities for its event projectors, and L\"uders conditioning on that same
commuting algebra.

This regulated architecture should be read as a reference implementation, not as the unique OPH
UV completion. The microphysics paper is explicit on that boundary: it selects one finite
gauge-register or quantum-link style screen model because it is local, finite-dimensional,
gauge-aware, and naturally compatible with collars, overlap observables, and repair maps. The
particle paper inherits that architecture as a concrete screen-side realization of the OPH
language, while leaving uniqueness and continuum lift open.

\subsection{External inputs and where they enter the particle derivation}

The quantitative derivation developed here uses two external continuous inputs:
\begin{align}
P &\equiv a_{\mathrm{cell}}/\ell_P^2 = 1.63094,\\
N_{\mathrm{scr}} &\equiv \log \dim \mathcal H_{\mathrm{tot}} \sim 10^{122}.
\end{align}

The particle-spectrum draft needs to keep these inputs visible because several later sections
depend on them in different ways.

\paragraph{Pixel area \(P\).}
The pixel area is the declared local UV-area input. In this derivation it feeds the
heat-kernel or edge-law calibration stage and, through that route, the electroweak
closure surface. It is the upstream numerical
input for the electroweak, Higgs/top, flavor, quark, and hadron-facing branches. The particle
paper therefore treats \(P\) as a declared input to the quantitative
derivation, not as a quantity derived inside the particle-spectrum argument itself.

\paragraph{Screen capacity \(N_{\mathrm{scr}}\).}
The screen capacity is not part of the local recovered-core gravity/gauge derivation. It enters
the separate screen-capacity branch once the screen-capacity identification is adopted. In the
particle context, this matters mainly for the capacity-based neutrino estimate and for the
cosmological-capacity discussion that frames why local null data do not determine the
cosmological constant. The particle paper should therefore keep the neutrino chapter honest:
capacity-based order-of-magnitude information and the neutrino derivation developed here are not the
same thing.

\paragraph{Why these are the only declared continuous inputs.}
One of the discipline constraints of the OPH derivation is that particle-sector freedom should not
be hidden in a large uncontrolled parameter set. This paper compresses the
continuous external freedom into \(P\) and \(N_{\mathrm{scr}}\), then asks the overlap,
modular, gauge, and admissibility machinery to do the rest. Whether later continuation branches
remain faithful to that discipline is one of the central questions this draft paper is meant to
make explicit. For the particle-spectrum branch specifically, the nontrivial common calibration
burden sits on \(P\): \(N_{\mathrm{scr}}\) enters the separate capacity/cosmology side, whereas
the reverse-engineering claim for masses and couplings is that one shared pixel constant should
drive the spectrum instead of a long sector-by-sector input list.

\subsection{Technical premises carried into the particle-spectrum derivation}

The particle paper uses the same technical-premise boundary as
\emph{Recovering Relativity and Standard Model Structure from Observer Overlap Consistency}~\cite{ophcompact}, together with
the regulator-language clarification of the technical supplement. The important point is not to
memorize every label, but to separate what is branch-internal from what is still theorem-external.

The third axiom internalizes several pieces of infrastructure that older drafts often
carried as separate assumptions: the local Gibbs form of the regulator-scale state, the
quasi-local propagation bound, the endpoint-control estimate for bounded intervals, and the
meaning of refinement stability itself. In regulator language, one works with finite local
Hilbert spaces and a boundary-fixed overlap action on cut data, often summarized as the R0/R1
presentation. On top of that, the still-genuine extra burdens are the controlled scaling-limit
scope, the geometric modular branch, fixed-cap stationarity in the Jacobson step, the bosonic
and fiber-functor clauses required by compact-gauge reconstruction, and the vanishing gluing
obstruction when global transportability is invoked.
For the BW/geometric side, the smallest honest internal extension route is a two-step chain:
first a scaling-limit cap-pair extraction theorem that produces transported cap limits with
vanishing carried collar errors, and then an ordered null cut-pair rigidity theorem that
collapses the residual cap-preserving conformal freedom on those extracted cap pairs to the
hyperbolic subgroup.

For the particle-spectrum derivation, this has a practical consequence. Structural statements such
as compact-gauge reconstruction, the realized Standard Model quotient, the exact hypercharge
lattice, the generation count \(N_g=3\), the color count \(N_c=3\), and the symmetry-protected
massless carrier zeros are not floating independently of the OPH premise map. They live on a
controlled chain whose scope conditions need to remain visible inside the particle paper,
especially once later chapters turn from structural branches to calibration, continuation, and
nonperturbative sectors.

\subsection{Core theorem packages used later}

The particle paper will repeatedly use earlier OPH results without re-proving them in full. The
theorem packages imported into the present draft are the following.

\begin{enumerate}[leftmargin=2em]
\item \textbf{Patch-net and overlap language.} The consensus paper organizes overlap repair,
fixed-point language, cycle holonomy, and gauge-as-gluing in a form that is especially useful
for flavor transport and observer-centric interpretation.
\item \textbf{Screen-side regulated architecture.} The screen-microphysics paper gives the
finite gauge-register screen model, the regulated patch-net embedding theorem, and the fixed-cutoff
edge heat-kernel / Casimir theorem.
\item \textbf{Measurement interface.} The same microphysics paper, together with the technical
supplement and the observer wrapper, supplies the fixed-cutoff central-record, Born-rule, and
L\"uders-conditioning package used when later chapters discuss how measurements pick out definite
observer-accessible outcomes.
\item \textbf{Compact gauge reconstruction and Standard Model structure.}
\emph{Recovering Relativity and Standard Model Structure from Observer Overlap Consistency}~\cite{ophcompact}, the
main paper, and the dedicated gauge-group fragment together provide the path from refinement-stable
edge sectors to a compact group, then from MAR to
\[
\frac{\SU(3)\times\SU(2)\times\U(1)}{\mathbb Z_6},
\qquad
N_g=3,
\qquad
N_c=3.
\]
\item \textbf{Continuation-level topic notes.} The technical supplement carries useful material
on supersymmetry, baryogenesis, proton stability, generation structure, and Yukawa hierarchy,
but those sections have to be read with their stated status boundaries intact. They are sources
for later discussion and continuation chapters, not automatic upgrades to recovered-core theorems.
\end{enumerate}

This reading rule will govern the whole paper. Structural gauge and carrier results can be used
as structural results. Electroweak calibration can be used as calibration. The
Higgs/top critical stage can be used as a secondary quantitative stage. The quark,
charged-lepton, neutrino, and hadron chapters retain their stated continuation or
simulation status unless the underlying derivation chains close more strongly.

\section{Gauge Architecture, Standard Model Structure, and Structural Carriers}

The particle-spectrum derivation does not start from a list of particles. It starts from the
compact-gauge branch. OPH first reconstructs a compact gauge
group from the refinement-stable transportable edge-sector category and only then applies MAR to
select the realized low-energy branch. That logical order matters for the particle paper. The
gauge branch fixes the structural carrier content before any mass readout, and it also explains
why some rows are exact structural zeros rather than fitted or accidentally small numbers.

On the realized branch, the gauge/content result used throughout this paper is
\[
\frac{\SU(3)\times\SU(2)\times\U(1)}{\mathbb Z_6},
\qquad
N_g=3,
\qquad
N_c=3.
\]
\emph{Recovering Relativity and Standard Model Structure from Observer Overlap Consistency}~\cite{ophcompact}, the longer main derivation surface, and the dedicated gauge-group proof
fragment all agree on the logical split: compact-gauge reconstruction first, realized Standard
Model selection second. The first stage gives a compact internal symmetry group under the stated
categorical premises. The second stage uses MAR, anomaly cancellation, the one-Higgs admissible
matter package, intrinsic CP capability, weak-sector UV completeness, and Witten-parity input to
select the realized quotient, the generation count, and the color count.

For the particle zoo this structural branch fixes three exact carrier statements.
First, the photon is massless on the realized electromagnetic branch. Second, the gluons are
massless on the color branch, even though free gluons are never seen as asymptotic particles
because color is confined. Third, the graviton is massless on the dynamical-metric branch of the
OPH gravity chain. These are not continuation-level masses but symmetry-protected zeros of the
realized branch.

\begin{table}[h]
\centering
\small
\begin{tabular}{@{}>{\raggedright\arraybackslash}p{0.13\textwidth}>{\raggedright\arraybackslash}p{0.12\textwidth}>{\raggedright\arraybackslash}p{0.12\textwidth}>{\raggedright\arraybackslash}p{0.23\textwidth}>{\raggedright\arraybackslash}p{0.26\textwidth}@{}}
\toprule
Carrier & Status & OPH output & Physical role & Why the mass is zero \\
\midrule
Photon & structural & \(0\) GeV & electromagnetic force carrier & gauge redundancy forbids a hard mass \\
Gluons & structural & \(0\) GeV & color-force carriers & color-gauge redundancy forbids a hard mass \\
Graviton & structural & \(0\) GeV & spin-2 metric carrier & diffeomorphism redundancy forbids a hard mass \\
\bottomrule
\end{tabular}
\end{table}

These structural zeros should not be confused with the later quantitative particle rows. The
particle paper should therefore treat the photon, gluons, and graviton as consequences of the
gauge and gravity branches themselves, not as outputs of the later calibration or flavor
machinery. The later boson sections will add the \(W\) boson, the \(Z\) boson, and the Higgs
boson, but those are not structural zeros; they sit on electroweak and Higgs/top quantitative stages.

\section{Electroweak Calibration and the Higgs/Top Critical Stage}

The reported bosonic sector distinguishes sharply between emitted numerical rows and
missing rows. The active bosonic stages are electroweak
calibration stage and the Higgs/top critical stage, while charged-lepton rows
and neutrino rows remain explicit gaps and hadron rows are suppressed. The bosonic numerical
sector is therefore concentrated in two places: electroweak calibration for the \(W\) boson and \(Z\) boson, and
the Higgs/top critical stage for the Higgs boson, together with the top quark.

\subsection{Electroweak calibration stage}

The electroweak derivation consists of a single-\(P\) running family followed by a
reduced two-scalar carrier, a selector on that carrier, an exact carrier mass chart, and
then a target-free source-only repair theorem beyond the present carrier. In practical terms, the construction starts from
the declared pixel input \(P\), builds the running electroweak family, reduces it to the present
carrier, reads the \(W\) boson/\(Z\) boson pair from the selected point on that carrier,
and then emits the coherent repaired electroweak quintet from the source basis alone.
Calibration here means exactly that ordering: first fix the shared pixel scale \(P\) on the
declared D10 running/matching surface, thereby fixing the source basis
\[
(\alpha_{2,m_Z},\alpha_{Y,m_Z},\eta_{\mathrm{source}},v),
\]
and only then emit the repaired electroweak quintet forward from that calibrated basis. The
active public \(W/Z\) rows are therefore not obtained by re-inserting the target masses after
calibration; they are forward outputs of the closed D10 calibration theorem.

The public reported rows are
\[
m_W = 80.37700001539531~\mathrm{GeV},
\qquad
m_Z = 91.18797807794321~\mathrm{GeV}.
\]
These numbers are emitted on the target-free source-only repair theorem, not on the unmodified
selected carrier. The exact current-carrier chart remains explicit on disk and emits
the local pair
\[
m_W^{\mathrm{carrier}} = 80.38629169244275~\mathrm{GeV},
\qquad
m_Z^{\mathrm{carrier}} = 91.18290444674243~\mathrm{GeV},
\]
so the mathematics distinguishes three electroweak objects: the exact current-carrier chart, the
freeze-once coherent repair surface, and the target-free source-only theorem surface that supplies
the public rows.

\begin{theorem}[Freeze-once coherent electroweak repair law]
Let the emitted running/core electroweak basis be
\[
Q_{\mathrm{run}}(P)=
\bigl(\alpha_{Y,m_Z}(P),\alpha_{2,m_Z}(P),v(P),\eta_{\mathrm{source}}(P)\bigr),
\]
and write
\[
\tau_Y(\tau_2)
=
-\frac{\tau_2+2\eta_{\mathrm{source}}}{1+4\tau_2^2},
\qquad
n_{\mathrm{fiber}}(\tau_2)
=
1+\frac{\alpha_{Y,m_Z}\tau_Y(\tau_2)+\alpha_{2,m_Z}\tau_2}
{\alpha_{Y,m_Z}+\alpha_{2,m_Z}}.
\]
Freeze one authoritative target pair
\[
\mathcal T^\dagger=(M_W^\dagger,M_Z^\dagger),
\qquad 0<M_W^\dagger<M_Z^\dagger,
\]
and define the charged anchor
\[
\tau_{2,W}^\dagger
=
\frac{M_W^{\dagger\,2}}{\pi v^2\alpha_{2,m_Z}}-1,
\qquad
\delta\alpha_2^\dagger
=
\frac{M_W^{\dagger\,2}}{\pi v^2}-\alpha_{2,m_Z}.
\]
Then the fiber-parallel hypercharge leg is
\[
\tau_Y^\dagger=\tau_Y(\tau_{2,W}^\dagger),
\qquad
n_{\mathrm{fiber}}^\dagger=n_{\mathrm{fiber}}(\tau_{2,W}^\dagger),
\]
\[
M_{Z,\mathrm{fiber}}^\dagger
=
v\sqrt{\pi(\alpha_{Y,m_Z}+\alpha_{2,m_Z})\,n_{\mathrm{fiber}}^\dagger},
\]
and the orthogonal neutral-shear scalar is
\[
\delta M_Z^\perp=M_Z^\dagger-M_{Z,\mathrm{fiber}}^\dagger,
\]
equivalently
\[
\delta n^\dagger
=
\frac{(M_Z^\dagger+M_{Z,\mathrm{fiber}}^\dagger)\,\delta M_Z^\perp}
{\pi v^2(\alpha_{Y,m_Z}+\alpha_{2,m_Z})},
\qquad
\delta\alpha_Y^\perp
=
\frac{(M_Z^\dagger+M_{Z,\mathrm{fiber}}^\dagger)\,\delta M_Z^\perp}{\pi v^2}.
\]
The fiber-parallel hypercharge motion is
\[
\delta\alpha_Y^\parallel
=
\alpha_{Y,m_Z}\,
\frac{8\eta_{\mathrm{source}}(\tau_{2,W}^\dagger)^2-\tau_{2,W}^\dagger}
{1+4(\tau_{2,W}^\dagger)^2}.
\]
Therefore the unique coherent repair package is
\[
\Sigma_{EW}^\dagger
=
\bigl(\delta\alpha_2^\dagger,\delta\alpha_Y^\parallel,\delta\alpha_Y^\perp\bigr),
\]
equivalently \((\tau_{2,W}^\dagger,\delta n^\dagger)\). If the selected carrier anchor is the
\(\tau_2=0\) fiber point, so that
\[
\alpha_{2,*}=\alpha_{2,m_Z},
\qquad
\alpha_{Y,*}=\alpha_{Y,m_Z}(1-2\eta_{\mathrm{source}}),
\]
then the repaired coherent couplings are
\[
\alpha_{2,\dagger}=\alpha_{2,*}+\delta\alpha_2^\dagger,
\qquad
\alpha_{Y,\dagger}=\alpha_{Y,*}+\delta\alpha_Y^\parallel+\delta\alpha_Y^\perp,
\]
and they emit one coherent quintet
\[
M_W^\dagger=v\sqrt{\pi\alpha_{2,\dagger}},
\qquad
M_Z^\dagger=v\sqrt{\pi(\alpha_{Y,\dagger}+\alpha_{2,\dagger})},
\]
\[
\alpha_{\mathrm{em},\dagger}^{-1}
=
\frac{\alpha_{Y,\dagger}+\alpha_{2,\dagger}}
{\alpha_{Y,\dagger}\alpha_{2,\dagger}},
\qquad
\sin^2\theta_{W,\dagger}
=
\frac{\alpha_{Y,\dagger}}{\alpha_{Y,\dagger}+\alpha_{2,\dagger}}.
\]
So the frozen-target electroweak repair is not another chart search or selector shell. It is one
unique coherent value-emission law on top of the already-closed fiber map.
\end{theorem}

For the current local PDG-2025 calibration surface,
\[
\alpha_{2,m_Z}=0.03377843630219015,
\qquad
\alpha_{Y,m_Z}=0.010131601067241624,
\qquad
\eta_{\mathrm{source}}=0.022147000871961295,
\]
\[
v=246.76711732749683~\mathrm{GeV},
\qquad
M_W^\dagger=80.377~\mathrm{GeV},
\qquad
M_Z^\dagger=91.18797809193725~\mathrm{GeV},
\]
and the law emits
\[
\tau_{2,W}^\dagger=-2.3116268316325517\times10^{-4},
\qquad
\delta\alpha_2^\dagger=-7.808313968675484\times10^{-6},
\]
\[
M_{Z,\mathrm{fiber}}^\dagger=91.17717018044121~\mathrm{GeV},
\qquad
\delta M_Z^\perp=10.807911496044653~\mathrm{MeV},
\]
\[
\delta\alpha_Y^\parallel=2.3421435088707883\times10^{-6},
\qquad
\delta\alpha_Y^\perp=1.0302892666790446\times10^{-5},
\qquad
\delta n^\dagger=2.3463639031113336\times10^{-4}.
\]
So the repaired coherent couplings are
\[
\alpha_{2,\dagger}=0.03377062798822147,
\qquad
\alpha_{Y,\dagger}=0.00969547694807616,
\]
and the same coherent branch emits
\[
M_W^\dagger=80.377~\mathrm{GeV},
\qquad
M_Z^\dagger=91.18797809193725~\mathrm{GeV},
\]
\[
\alpha_{\mathrm{em},\dagger}^{-1}=132.752408550967,
\qquad
\sin^2\theta_{W,\dagger}=0.223058333896849.
\]
That freeze-once law provides an exact coherent validation surface. The active public electroweak
surface is the target-free source-only repair law.

\begin{theorem}[Target-free source-only electroweak repair law]
Let the D10 source-only emitted basis be
\[
Q_{\mathrm{src}}(P)=
\bigl(\alpha_{2,m_Z}(P),\alpha_{Y,m_Z}(P),\eta_{\mathrm{source}}(P),v(P)\bigr),
\]
and define
\[
\beta_{\mathrm{EW}}
:=
\frac{\alpha_{2,m_Z}-\alpha_{Y,m_Z}}{\alpha_{2,m_Z}+\alpha_{Y,m_Z}},
\qquad
\alpha_u^{\mathrm{seed}}
:=
\frac{\eta_{\mathrm{source}}}{\beta_{\mathrm{EW}}},
\qquad
\lambda_{\mathrm{EW}}
:=
\frac{\eta_{\mathrm{source}}\alpha_u^{\mathrm{seed}}}{4}
=
\frac{\eta_{\mathrm{source}}^2}{4\beta_{\mathrm{EW}}}.
\]
Then the beyond-current-carrier repair is uniquely
\[
\tau_2^{\mathrm{exact}}
=
-\lambda_{\mathrm{EW}}
\left(
1+\frac23\eta_{\mathrm{source}}
+
\left(1-\frac{\beta_{\mathrm{EW}}}{6}\right)\eta_{\mathrm{source}}^2
\right),
\]
\[
\delta n^{\mathrm{exact}}
=
\lambda_{\mathrm{EW}}
\left(
1+\frac43\eta_{\mathrm{source}}
+
\left(2-\frac{\beta_{\mathrm{EW}}}{6}\right)\eta_{\mathrm{source}}^2
\right),
\]
with fiber hypercharge transport
\[
\tau_Y^{\mathrm{fiber}}
=
-\frac{\tau_2^{\mathrm{exact}}+2\eta_{\mathrm{source}}}
{1+4(\tau_2^{\mathrm{exact}})^2}.
\]
The repaired coupling shifts are
\[
\delta\alpha_2=\alpha_{2,m_Z}\tau_2^{\mathrm{exact}},
\qquad
\delta\alpha_Y^{\parallel}
=
\alpha_{Y,m_Z}
\frac{8\eta_{\mathrm{source}}(\tau_2^{\mathrm{exact}})^2-\tau_2^{\mathrm{exact}}}
{1+4(\tau_2^{\mathrm{exact}})^2},
\]
\[
\delta\alpha_Y^{\perp}
=
(\alpha_{2,m_Z}+\alpha_{Y,m_Z})\,\delta n^{\mathrm{exact}}.
\]
With
\[
\alpha_{2,*}=\alpha_{2,m_Z},
\qquad
\alpha_{Y,*}=\alpha_{Y,m_Z}(1-2\eta_{\mathrm{source}}),
\]
the repaired couplings are
\[
\alpha_2'=\alpha_{2,*}+\delta\alpha_2,
\qquad
\alpha_Y'=\alpha_{Y,*}+\delta\alpha_Y^{\parallel}+\delta\alpha_Y^{\perp},
\]
and they emit one coherent family
\[
W=v\sqrt{\pi\alpha_2'},
\qquad
Z=v\sqrt{\pi(\alpha_2'+\alpha_Y')},
\]
\[
\alpha_{\mathrm{em}}^{-1}
=
\frac{\alpha_2'+\alpha_Y'}{\alpha_2'\alpha_Y'},
\qquad
\sin^2\theta_W
=
\frac{\alpha_Y'}{\alpha_2'+\alpha_Y'},
\qquad
v=v(P).
\]
So the electroweak repair is emitted from the D10 source basis alone, with no frozen target input,
no inverse slice, and no mixed readout.
\end{theorem}

Using the current D10 source basis
\[
\alpha_{2,m_Z}=0.03377843630219015,
\qquad
\alpha_{Y,m_Z}=0.010131601067241624,
\qquad
\eta_{\mathrm{source}}=0.022147000871961295,
\]
\[
v=246.76711732749683~\mathrm{GeV},
\]
the theorem emits
\[
W=80.37700001539531~\mathrm{GeV},
\qquad
Z=91.18797807794321~\mathrm{GeV},
\]
\[
\alpha_{\mathrm{em}}^{-1}=132.7524088191668,
\qquad
\sin^2\theta_W=0.22305833336075578.
\]
The freeze-once coherent repair surface provides compare-only validation, with machine-scale deltas
\[
\Delta M_W\approx 1.54\times 10^{-8}~\mathrm{GeV},
\qquad
\Delta M_Z\approx -1.40\times 10^{-8}~\mathrm{GeV}.
\]

The electroweak calibration branch also contains three subordinate theorem objects: an
underdetermination theorem for the quadratic repair family, the smallest conditional route
through \(\mathrm{ColorBalancedQuadraticRepairDescent}_{EW}\), and the source-only expression
\(\mathrm{EWTargetEmitter}_{D10}\) whose promotion yields the active theorem. These objects locate
the public repair law inside the calibration branch without changing the emitted public surface.

\subsection{Higgs/top critical stage}

The Higgs/top critical stage is a secondary quantitative stage built on top of the electroweak gauge core.
It is described by a compact shared seed together with a Jacobian readout map. The reported
rows use a promoted forward seed whose certificate records the exact fixed-ray forward
path. This gives
\[
m_H = 125.218922~\mathrm{GeV},
\qquad
m_t = 172.388646~\mathrm{GeV}.
\]
The Higgs boson belongs naturally in the boson discussion, while the top quark will be revisited
in the quark-family chapter as the third-generation up-type quark. The important point here is
stage logic: the reported top-quark row is carried by the Higgs/top critical stage, not by the quark derivation itself.

The important distinction is between the closed forward seed and the broader paper surface. The
forward seed is closed on its emitted one-scalar branch, so the Higgs/top rows are
not just loose candidates. What remains open is presentation-level rigor and synchronization
across the broader OPH exposition, not the existence of a larger missing mass-moving
coordinate on this branch.

\begin{table}[h]
\centering
\small
\begin{tabular}{@{}llll@{}}
\toprule
Particle & Stage & Status & Reported value \\
\midrule
\(W\) boson & electroweak calibration & calibration & \(80.37700001539531~\mathrm{GeV}\) \\
\(Z\) boson & electroweak calibration & calibration & \(91.18797807794321~\mathrm{GeV}\) \\
Higgs boson & Higgs/top critical stage & secondary quantitative & \(125.218922~\mathrm{GeV}\) \\
Top quark & Higgs/top critical stage & secondary quantitative & \(172.388646~\mathrm{GeV}\) \\
\bottomrule
\end{tabular}
\end{table}

The electroweak chapter therefore presents three mathematically distinct surfaces: the reported
target-free source-only theorem rows, the current-carrier chart beneath them, and the freeze-once
coherent repair surface used for compare-only validation. That is the technically honest way to
present the bosonic sector.

\section{Flavor Transport, Generation Structure, and the Yukawa Dictionary}

The gauge branch explains why the realized low-energy sector has three generations, but it does
not yet by itself produce the full flavor dictionary. The particle paper therefore needs a
separate flavor chapter. This chapter is the bridge between the structural result \(N_g=3\) and
the matter-family chapters for quarks, charged leptons, and neutrinos.

The flavor derivation is deliberately constructive. It does not claim that the full
OPH flavor observable is theorem-level. Instead, it builds the object chain that a
closed theorem would have to pass through. The derivation chain has the following main
architecture: start from a refinement-indexed family transport kernel, derive a centered
generation-bundle branch generator, lift that to same-label transport data, derive the induced
overlap-edge transport cocycle, reduce the cocycle to a persistent flavor observable, and then
push that observable into common sector-response objects for the downstream quark, charged-lepton,
and neutrino derivations.

Mathematically, this derivation is where ``why this family exists'' becomes a concrete technical
question. The realized Standard Model branch gives three generations structurally; the flavor
derivation is where one tries to turn that three-generation existence statement into transport,
splitting, suppression, phase, and excitation data. The relevant objects are not yet the final
fermion masses themselves but the intermediate transport and spectral structures that later mass
readouts consume.

The active chain, in slightly compressed form, is:
\begin{enumerate}[leftmargin=2em]
\item normalize a refinement-indexed family transport kernel;
\item derive a centered generation-bundle branch generator on the realized three-generation
charged bundle;
\item lift this to same-label edge-line transport and then to an overlap-edge cocycle with
explicit defect and gap bookkeeping;
\item reduce those data to projectors, spectral gaps, pair suppressions, and cycle phases;
\item push the resulting family object into sector-response objects and then into
suppression/phase tensors for the downstream matter derivations.
\end{enumerate}

Three continuation problems are attached to this region of the derivation. The
first is the Yukawa/excitation dictionary itself. The second is a microphysical bridge that
tries to carry edge statistics into that dictionary if the direct theorem route stalls. The third
is the quark-texture / CKM / CP problem. Together they mark the point where the flavor derivation
stops at constructive intermediate objects rather than closed particle data.

\section{Quark Family Derivation}

Among the matter families, quarks are the only family with reported emitted numerical rows. But
even here the status is not ``closed theorem.'' The reported quark rows are all
continuation-level rows, and the reason is explicit: the quark derivation
implements the local mean split, the quark descent, the forward Yukawa surface, and the
pure-\(B\) source-readback shell, but the light-quark splitting remains on an overlap
continuation frontier rather than on the recovered-core tier.

\subsection{What quarks are in this derivation}

Quarks are the color-charged elementary constituents of hadronic matter. In nature, the up quark
and down quark dominate protons and neutrons, while the strange quark, charm quark, bottom
quark, and top quark appear in progressively heavier and more unstable sectors. In the OPH
particle derivation, quarks are the first matter family for which the flavor chain
emits mass-facing reported candidates rather than only gap surfaces.

The quark path takes the shared flavor data, emits the quark-sector mean split, assembles
the quark descent, builds the forward Yukawa matrices, and fixes a pure-\(B\) source-readback
shell. Along that derivation, the unresolved mass-moving object is the emitted
pure-\(B\) up/down readback payload pair, namely the up-side readback
\texttt{source\_readback\_u\_log\_per\_side} together with the down-side readback
\texttt{source\_readback\_d\_log\_per\_side}. From that pair the odd amplitudes, the \(\tau\)
pair, and the final odd diagonal contribution are then read back algebraically.

The strongest tier conclusion on top of that builder view is that the available OPH corpus does
not support a nonzero light-quark pure-\(B\) source selector at
recovered-core tier. So the broader honest open object is not ``find one more recovered-core
quark scalar.'' It is: open a light-quark isospin-breaking selector / overlap-defect scalar
branch without pretending that the light \(u/d\) split is part of the recovered core.

\subsection{Current emitted quark rows}

\begin{table}[h]
\centering
\small
\begin{tabular}{@{}llll@{}}
\toprule
Quark & Status & Reported value & Derivation stage \\
\midrule
Up quark & continuation & \(0.002287663517~\mathrm{GeV}\) & overlap continuation branch \\
Down quark & continuation & \(0.005121438948~\mathrm{GeV}\) & overlap continuation branch \\
Strange quark & continuation & \(0.09161114241~\mathrm{GeV}\) & overlap continuation branch \\
Charm quark & continuation & \(1.25669206~\mathrm{GeV}\) & overlap continuation branch \\
Bottom quark & continuation & \(3.82011788~\mathrm{GeV}\) & overlap continuation branch \\
Top quark & secondary quantitative & \(172.388646~\mathrm{GeV}\) & Higgs/top critical stage \\
\bottomrule
\end{tabular}
\end{table}

These comparisons are made against running quark references, not direct free-particle pole
masses, because quarks are confined. That is not a technicality; it is part of what the quark
chapter must explain to technical readers. The reported quark rows should therefore be interpreted
as OPH readouts against the standard running-mass comparison surface, not as
asymptotic free-particle masses.

\subsection{Strongest overlap-branch mass-side and mixing-side primitives}

The strongest constructive continuation object is no longer a one-scalar light-sector overlap law.
On the ordered-family overlap-continuation branch, the quark continuation splits cleanly into a
two-scalar mass side and a matrix-valued mixing side.

Write
\[
\Delta := \Delta_{ud}^{\mathrm{overlap}},
\qquad
\Lambda := \Lambda_{ud}^{B,\mathrm{transport}}
=
\frac{\sigma_u\sigma_d}{2(\sigma_u+\sigma_d)}\,\Delta,
\]
\[
\tau_u
=
\frac{\Lambda}{\sigma_u}
=
\frac{\Delta}{2}\frac{\sigma_d}{\sigma_u+\sigma_d},
\qquad
\tau_d
=
\frac{\Lambda}{\sigma_d}
=
\frac{\Delta}{2}\frac{\sigma_u}{\sigma_u+\sigma_d}.
\]
On the emitted spread package,
\[
\sigma_u = 5.5905,
\qquad
\sigma_d = 3.3049.
\]
The odd transport therefore collapses to one overlap scalar \(\Delta\). The even transport
collapses independently to one centered quadratic scalar
\[
\eta_Q^{\mathrm{centered}},
\qquad
Q_{\mathrm{ord}}
=
\frac{1-x_2^2}{3}(1,-2,1),
\qquad
x_2=-0.5175863354681689,
\]
with equivalent parameter
\[
\kappa_Q
=
\frac{\eta_Q^{\mathrm{centered}}}{2(1-x_2^2)},
\qquad
E_{\mathrm{even}}
=
\frac{\eta_Q^{\mathrm{centered}}}{6}(1,-2,1).
\]
So the overlap-continuation mass side reduces to the pair
\[
(\Delta_{ud}^{\mathrm{overlap}},\,\eta_Q^{\mathrm{centered}}).
\]

There is also a sharper same-family specialization beneath that two-scalar description. Let
\[
C_1(Y)=\frac35Y^2,
\qquad
Y(u^c)=-\frac23,
\qquad
Y(d^c)=\frac13.
\]
Then the same-family overlap selector law gives
\[
\Delta_{ud}^{\mathrm{overlap}}
=
t_1\bigl(C_1(u^c)-C_1(d^c)\bigr)
=
\frac{t_1}{5},
\]
and on the compact same-family branch one also has
\[
\eta_Q^{\mathrm{centered}}
=
-\frac{1-x_2^2}{27}\,t_1.
\]
Eliminating \(t_1\) gives the exact branch law
\[
\eta_Q^{\mathrm{centered}}
=
-\frac{5(1-x_2^2)}{27}\,\Delta_{ud}^{\mathrm{overlap}}
\approx
-0.1355748861734504\,\Delta_{ud}^{\mathrm{overlap}}.
\]
So the compact same-family continuation does not honestly emit a distinguished scalar \(t_1\). It
emits only the one-parameter ray
\[
\mathcal R_{D12}^{ud}
:=
\left\{
\lambda\left(\frac15,-\frac{1-x_2^2}{27}\right):\lambda\ge 0
\right\},
\]
equivalently the one-parameter forward family \(\mathcal M_{D12}^{ud}(\lambda)\). Any numerical
\(t_1\) is therefore only a valid choice of \(\lambda\) on this ray, not a theorem-grade branch
emission.

If one specializes to the current exact-mean compare surface, the best point is
\[
\alpha_u = 1.0007763698011345,
\qquad
\alpha_d = 1.008463281557513,
\]
\[
\Delta_{ud}^{\mathrm{overlap}} = 0.14049991320632976,
\qquad
\eta_Q^{\mathrm{centered}} = -0.018104730181494357,
\]
with
\[
(u,c,t)
=
(0.0021599234039183404,\ 1.272939334772138,\ 172.36668155105923)\ \mathrm{GeV},
\]
\[
(d,s,b)
=
(0.004700052692857295,\ 0.09350445861546268,\ 4.182753646232867)\ \mathrm{GeV},
\]
and RMS log-mass error
\[
5.2130513964914564\times10^{-5}.
\]
Those values are not promoted as OPH-forward laws; they are the strongest current diagnostic
specialization of the two-scalar continuation family.

If one instead keeps the retained same-family sample point
\[
t_1^{\mathrm{sample}}=\lambda^{\mathrm{sample}}=0.6695617711471163,
\]
then
\[
\Delta_{ud}^{\mathrm{overlap}}=0.13391235422942327,
\qquad
\eta_Q^{\mathrm{centered}}=-0.018155152181872827,
\]
with
\[
(u,c,t)
=
(0.002176632492997213,\ 1.2566921714397312,\ 172.85192993931355)\ \mathrm{GeV},
\]
\[
(d,s,b)
=
(0.004708229528844486,\ 0.09160827127333702,\ 4.155513989984608)\ \mathrm{GeV},
\]
and RMS log error
\[
0.01078398352557181.
\]
These are sample forward masses on the emitted ray, not theorem-emitted masses of a
branch-selected \(\lambda^\star\).

The mixing side is closed on the same D12 continuation branch. Once the forward Yukawa step emits
the ordered same-label left eigenframes \(U_u\) and \(U_d\), the canonical transport unitary is
\[
V_{\mathrm{CKM}}^{\mathrm{fwd}} = U_u^\dagger U_d.
\]
After a diagonal rephasing gauge choice that makes \(V_{ud},V_{us},V_{cs},V_{cb},V_{tb}\)
real-positive, the principal unitary logarithm exists uniquely, and after one further diagonal
conjugation the generator takes the gauge-fixed form
\[
K_{\mathrm{CKM}}^{\mathrm{gf}}
=
i\,\mathrm{diag}(\chi_1,\chi_2,\chi_3)
+ \theta_{12}^{(K)}(E_{12}-E_{21})
+ \theta_{23}^{(K)}(E_{23}-E_{32})
+ \theta_{13}^{(K)}\bigl(e^{i\phi_K}E_{13}-e^{-i\phi_K}E_{31}\bigr).
\]
On the current retained same-family sample point this yields the invariant package
\[
\bigl(\theta_{12}^{(K)},\theta_{23}^{(K)},\theta_{13}^{(K)},\phi_K\bigr)
=
\bigl(0.007599799911069969,\ 0.001287397802967049,\ 3.0665238345096405\times 10^{-5},\ -1.85995254938665\bigr),
\]
\[
\chi =
\bigl(9.58567591\times 10^{-11},\ -1.91714078\times 10^{-10},\ 9.58570190\times 10^{-11}\bigr).
\]
The standard CKM angles and phase of the emitted matrix are
\[
\bigl(\theta_{12}^{\mathrm{CKM}},\theta_{23}^{\mathrm{CKM}},\theta_{13}^{\mathrm{CKM}},\delta_{\mathrm{CKM}}\bigr)
=
\bigl(0.00759980344251945,\ 0.0012874186416783875,\ 2.9643208230524648\times 10^{-5},\ 1.7011978493658415\bigr),
\]
\[
J = 2.8755912280868794\times 10^{-10},
\qquad
\|\exp(K_{\mathrm{CKM}}^{\mathrm{gf}})-V_{\mathrm{CKM}}^{\mathrm{fwd}}\|_F
\approx 5.6\times 10^{-16}.
\]
Equivalently, the same transport can be written as
\[
K_{\mathrm{CKM}}^{\mathrm{gf}}
=
i\,\mathrm{diag}(\chi_1,\chi_2,\chi_3)
+ \theta_{12}^{(K)}(E_{12}-E_{21})
+ \theta_{23}^{(K)}(E_{23}-E_{32})
+ \theta_{13}^{(K)}\bigl(e^{i\phi_K}E_{13}-e^{-i\phi_K}E_{31}\bigr).
\]
So the honest quark frontier is split in two:
\begin{enumerate}[leftmargin=2em]
\item one discrete left-handed same-label relative-sheet selector \(\sigma_{ud}\) that leaves the current D12 no-go class and changes the CKM moduli at all; and
\item after that branch choice, one mass-side scale-fixing law on the selected branch, independent of target masses and independent of CKM/CP.
\end{enumerate}

The smallest already-local finite basis orbit is diagnostic-only. Replacing the ordered same-label
left singular bases by right or conjugate-right singular bases can move the CKM angles, but those
chirality swaps leave the physical left-handed domain on which
\(V_{\mathrm{CKM}}^{\mathrm{fwd}}=U_u^\dagger U_d\) is defined. They therefore sharpen the
frontier by exclusion only: the remaining orbit must be left-handed and same-label.

\subsection{First generation: up quark and down quark}

The up quark and down quark are the lightest quarks and the basic constituents of protons and
neutrons. Both rows are reported numerically, but both remain continuation-level rows because the
odd \(B\)-mode completion is not yet closed.

\subsection{Second generation: strange quark and charm quark}

The strange quark and charm quark are heavier and mainly visible through unstable hadrons. Their
presence in the reported quark values is important because it shows that the OPH quark
derivation is not merely reproducing one light doublet; it propagates the family structure far
enough to emit second-generation rows on the same forward surface. But again, the paper
must resist overstating what this means. The existence of those rows does not mean the present D12
sheet is the physical quark branch, and it does not mean theorem-level odd-source completion has
been achieved.

\subsection{Third generation: bottom quark and top quark}

The bottom quark sits on the quark continuation surface. The top quark does not. The top
row is emitted by the Higgs/top critical stage, so the third generation is split across two quantitative
surfaces: a continuation-level bottom-quark row from the quark Yukawa branch and a
secondary-quantitative top-quark row from the Higgs/top critical stage. This mixed status
should remain explicit because it is one of the clearest examples of the mathematical
nonuniformity of the particle derivation across the particle zoo.

\subsection{What remains open in the quark sector}

Three open fronts remain visible and should be written exactly that way in the paper:
\begin{enumerate}[leftmargin=2em]
\item the Yukawa/excitation dictionary is not yet closed;
\item the builder-facing pure-\(B\) payload pair beneath the reported quark values is still open;
\item the broader honest overlap-continuation frontier is a wrong-branch no-go on the current D12 sheet, so the exact next object is one discrete relative-sheet selector \(\sigma_{ud}\); only after that branch selection does the mass-side ray \(\mathcal R_{D12}^{ud}\) ask for its own intrinsic scale law.
\end{enumerate}

So the correct headline for quarks is not ``fully derived.'' It is: quarks are the only matter
family with reported numerical rows, but the row-emission branch still sits above visible open
constructive objects.

\section{Charged-Lepton Family Derivation}

The charged-lepton derivation is one step less mature than the quark derivation. It contains an
ordered package, a support-obstruction certificate, a support-extension emitter shell, a
two-scalar completion-law shell, the smaller eta source-readback primitive beneath that shell,
and a forward shape/scale surface. But unlike the quark derivation, it does not yet emit
numerical masses. The electron, muon, and tau rows are therefore reported as \(n/a\).

The mathematical reason is clear. The charged-lepton derivation starts from
the ordered charged package, proves that the realized support is a one-dimensional linear subray,
exposes the canonical quadratic support-extension direction, maps that into the charged
excitation gaps, closes a two-scalar support-extension law shell, isolates the smaller eta
source-readback primitive on that same carrier, and then builds the log-spectrum and forward
shape/scale surface. The unresolved scalar order is
\[
\begin{aligned}
&\texttt{eta\_source\_support\_extension\_log\_per\_side},\\
&\texttt{sigma\_source\_support\_extension\_total\_log\_per\_side}.
\end{aligned}
\]
Until those objects are properly closed, the charged-lepton chapter remains a
continuation chapter rather than a quantitative one.

Physically, the electron is the light charged lepton that makes atoms and chemistry possible, the
muon is its heavier unstable cousin, and the tau lepton is the heaviest charged lepton. The role
of the charged-lepton chapter is to explain how those three rows emerge from the shared flavor
dictionary without falling back to Koide-assisted fitting. The exact burden is sharper than a
generic ``emit \(\eta\) and \(\sigma\)'' instruction: the theorem-grade lane still has to derive
promotion of the latent charged sector-response candidate
\(\widehat C_e^{\mathrm{cand}}\) to theorem-grade \(\widehat C_e\) by closing the branch-generator
splitting theorem, and then one theorem-grade affine-covariant absolute anchor
\(A_{\mathrm{ch}}\) for the charged common-shift quotient.

\begin{theorem}[Charged same-carrier source-pair readback]
Fix the ordered charged carrier
\[
(-1,x_2,1),
\qquad
x_2=-0.5175863354681689.
\]
If the charged source pair
\[
(\eta_{\mathrm{ext}},\sigma_{\mathrm{ext}})
=
(\eta_{\mathrm{source\_support\_extension\_log\_per\_side}},
\sigma_{\mathrm{source\_support\_extension\_total\_log\_per\_side}})
\]
is emitted on that carrier, then the centered charged logs are
\[
e_{\log,\mathrm{centered}}
=
-\frac{(3+x_2)\sigma_{\mathrm{ext}}-\eta_{\mathrm{ext}}}{6},
\]
\[
\mu_{\log,\mathrm{centered}}
=
\frac{x_2\sigma_{\mathrm{ext}}-\eta_{\mathrm{ext}}}{3},
\]
\[
\tau_{\log,\mathrm{centered}}
=
\frac{(3-x_2)\sigma_{\mathrm{ext}}+\eta_{\mathrm{ext}}}{6},
\]
and therefore the charged masses are
\[
m_e = g_e\,e^{e_{\log,\mathrm{centered}}},
\qquad
m_\mu = g_e\,e^{\mu_{\log,\mathrm{centered}}},
\qquad
m_\tau = g_e\,e^{\tau_{\log,\mathrm{centered}}}.
\]
\end{theorem}

This theorem is exact on the same-carrier shell, but the values themselves are still open.
The visible scalar order remains
\[
\eta_{\mathrm{ext}}
\quad\text{then}\quad
\sigma_{\mathrm{ext}},
\]
and the charged absolute scale \(g_e\) remains unresolved. Builder-facing code therefore still
exposes \(\eta_{\mathrm{ext}}\) and then \(\sigma_{\mathrm{ext}}\) as the first same-carrier
residuals, but the theorem-grade waiting set is sharper than that surface suggests. If the latent
candidate \(\widehat C_e^{\mathrm{cand}}\) is promoted, then \(\eta_{\mathrm{ext}}\) and
\(\sigma_{\mathrm{ext}}\) become charged spectral invariants rather than separate primitive goals,
and the absolute-scale burden is pushed to one affine-covariant absolute charged anchor
\(A_{\mathrm{ch}}\). In the
current local chain, \(\widehat C_e\) itself is undeclared: only the centered compressed
generation-bundle branch-operator candidate \(\widehat C_e^{\mathrm{cand}}\) is on disk, and the
remaining burden for that package is the upstream promotion theorem
\texttt{oph\_generation\_bundle\_branch\_generator\_splitting}, not a new ad hoc charged
operator ansatz. Its smallest remaining clause is the compression-descendant commutator statement
\texttt{compression\_descendant\_commutator\_vanishes\_or\_is\_uniformly\_quadratic\_small\_after\_central\_split}.

\begin{theorem}[Charged absolute-scale underdetermination]
Let
\[
E_e^{\mathrm{centered}}
=
\bigl(e_{\log,\mathrm{centered}},\mu_{\log,\mathrm{centered}},\tau_{\log,\mathrm{centered}}\bigr)
\]
be the centered charged log triple emitted from the charged source pair
\((\eta_{\mathrm{ext}},\sigma_{\mathrm{ext}})\). Then for every \(c\in\mathbb R\),
\[
Y_e(c):=\exp(c)\,\mathrm{diag}\!\bigl(e^{E_e^{\mathrm{centered}}}\bigr)
\]
has the same charged spectral invariants
\[
\eta_{\mathrm{ext}},\qquad \sigma_{\mathrm{ext}},\qquad \gamma_{21},\qquad \gamma_{32},
\]
the same centered log vector, and the same ratio data. Only the absolute masses scale:
\[
(m_e,m_\mu,m_\tau)\longmapsto e^c\,(m_e,m_\mu,m_\tau).
\]
Hence the present charged OPH chain determines only the quotient class
\[
E_e^{\mathrm{centered}}\in \mathbb R^3/\langle(1,1,1)\rangle,
\]
not the absolute scalar \(g_e=e^c\).
\end{theorem}

The proof is immediate from the centered sum rule
\[
e_{\log,\mathrm{centered}}+\mu_{\log,\mathrm{centered}}+\tau_{\log,\mathrm{centered}}=0,
\]
which implies
\[
\det Y_e^{\mathrm{shape}}=1,
\qquad
\det Y_e = g_e^3.
\]
All emitted charged invariants depend only on differences of log entries, so a common shift in the
\((1,1,1)\) direction leaves them unchanged. This is the exact reason the charged lane remains open:
the present theorem surface fixes the centered shape, not the absolute normalization.

One important code-level and mathematical cleanup is already complete: the charged absolute-scale
lane is explicitly typed. The charged scale is a linear quantity
\[
g_e = e^{\mu_e^{\mathrm{abs}}},
\]
so the log-coordinate \(\mu_e^{\mathrm{abs}}\) must not be mixed directly with centered log gaps.
The current-family writeback therefore records the type-consistent shell
\[
\mu_{e,\mathrm{seed}}^{\mathrm{abs}} = \log(0.9231656602589082) = -0.07994658034676537,
\]
\[
\mu_{e,\mathrm{cand}}^{\mathrm{abs}}
=
\mu_{e,\mathrm{seed}}^{\mathrm{abs}}-\gamma_{\min}
=
-0.38231224060567365,
\qquad
g_e^{\mathrm{cand}} = e^{\mu_{e,\mathrm{cand}}^{\mathrm{abs}}}=0.6822819838027987.
\]
This is a representation-consistency shell only, not a charged-mass theorem. It fixes the
linear-vs-log coordinate discipline while leaving the true emitted value law for \(g_e\) open.
The audited shortcut
\[
\Delta_e^{\mathrm{abs}}=0.30236566025890826,
\qquad
g_e=0.6822819838027987
\]
is not a theorem-grade closure. It merely chooses one representative on the common-shift orbit and
does not land on the physical charged masses.

The smallest honest future theorem object beyond this shell is therefore not an equalizer
already latent in the current surface, and not another ad hoc scalar fit for \(g_e\), but one
theorem-grade affine-covariant section \(A_{\mathrm{ch}}\) of the quotient map, satisfying
\[
A_{\mathrm{ch}}(\ell + c\,\mathbf 1)=A_{\mathrm{ch}}(\ell)+c.
\]
The clean realization of that section is an uncentered charged response lift carrying a
determinant line, in which
\[
A_{\mathrm{ch}}=\tfrac13 \log \det(Y_e)=\tfrac13 \operatorname{tr}(\log Y_e).
\]
Once such an anchor exists,
\[
g_e = e^{A_{\mathrm{ch}}(\ell)},
\qquad
\Delta_e^{\mathrm{abs}} = \log g_{\mathrm{ch}}^{\mathrm{shared}} - A_{\mathrm{ch}}(\ell).
\]

\paragraph{Strongest continuation bridge.}
There is, however, a sharper continuation-level bridge than the paper had before.  Assume:
\begin{enumerate}[leftmargin=2em]
\item a uniform \(\mathbb Z_6\) center-label ensemble, so \(\varepsilon = 1/6\);
\item the balanced Hermitian circulant branch, so \(|b|/a = 1/\sqrt2\);
\item the OPH phase choice \(\delta = 2/9\).
\end{enumerate}
With the scale-free normalization \(a=1\), the ordered roots
\[
r_k = a + 2|b|\cos\!\left(\delta + \frac{2\pi k}{3}\right)
\]
become
\[
r_e = 0.040349908219207475,
\qquad
r_\mu = 0.5802119201475368,
\qquad
r_\tau = 2.3794381716332555.
\]
This continuation bridge emits the same-carrier pair
\[
\eta_{\mathrm{ext}} = -6.729586682888832,
\qquad
\sigma_{\mathrm{ext}} = 8.154061112725994,
\]
with ordered gap values
\[
\gamma_{21}=5.33160859254774,
\qquad
\gamma_{32}=2.822452520178255,
\qquad
\kappa_{\mathrm{ext}}=-4.59605680397234,
\]
and centered logs
\[
E_{\log,\mathrm{centered}}
=
[-4.495223235091244,\ 0.836385357456495,\ 3.6588378776347503].
\]
Against the charged references, this continuation-centered shape has residual norm
\[
\left\|E_{\log,\mathrm{centered}}^{\mathrm{cont}}
-E_{\log,\mathrm{centered}}^{\mathrm{ref}}\right\|
\approx
2.13\times10^{-5},
\]
so the continuation branch is a near-exact shape closure up to one common absolute
scale. The compare-only common scale required for exact absolute masses is
\[
g_e^\star = 0.04577885783568762.
\]
Equivalently, relative to the stored shared-budget seed
\[
g_{\mathrm{ch}}^{\mathrm{shared}} = 0.9231656602589082,
\]
the missing determinant-normalization transport scalar would have to take the compare-only value
\[
\Delta_e^{\mathrm{abs},\star}
=
\log\!\frac{g_{\mathrm{ch}}^{\mathrm{shared}}}{g_e^\star}
=
3.003986333402356.
\]
This identifies the remaining unresolved objects sharply: the shape bridge is almost solved on
the continuation branch, but the theorem-grade derivation still lacks promotion of the latent
candidate \(\widehat C_e^{\mathrm{cand}}\) to theorem-grade \(\widehat C_e\), then one
affine-covariant absolute anchor \(A_{\mathrm{ch}}\), and therefore the emitted
\(\eta_{\mathrm{ext}}\), \(\sigma_{\mathrm{ext}}\), \(g_e\), and \(\Delta_e^{\mathrm{abs}}\)
that would follow.

\section{Neutrino Family Derivation}

The neutrino derivation is the strongest positive continuation result after the Higgs/top critical stage.
The flavor rows are still \(n/a\). The intrinsic isotropic branch is excluded exactly, the
weighted-cycle continuation branch lands in the observed PMNS/hierarchy regime, and one overall
positive normalization scalar remains open for the absolute neutrino scale.

The current weighted-cycle branch emits
\[
\theta_{12}=34.225904631810025^\circ,\qquad
\theta_{23}=49.72282845058266^\circ,\qquad
\theta_{13}=8.686355527700156^\circ,
\]
\[
\delta_{\mathrm{PMNS}}=305.58061231449796^\circ,\qquad
J=-0.02753115613565372,
\]
and the dimensionless hierarchy invariant
\[
\frac{\Delta m_{21}^2}{\Delta m_{32}^2}=0.030721110097966534.
\]
The exponent law is fixed by the positive transport-load segment between
\(\chi=1+\epsilon\) and \(1+\gamma_{1/2}\): on a one-dimensional affine segment,
the balanced selector and the least-distortion selector for any positive
translation-invariant quadratic form coincide at the midpoint, so
\[
D_\nu=\frac{\chi+(1+\gamma_{1/2})}{2},
\qquad
p=1+\gamma+\frac{\epsilon}{D_\nu}.
\]
With the compare-only atmospheric anchor
\[
\Delta m_{32}^2 = 2.438\times10^{-3}\ \mathrm{eV}^2,
\]
the same weighted-cycle branch gives
\[
(m_1,m_2,m_3)=(0.01745675648,\ 0.01948426065,\ 0.05308141307)\ \mathrm{eV},
\]
\[
\Delta m_{21}^2=7.489806641884242\times10^{-5}\ \mathrm{eV}^2,\qquad
\Delta m_{31}^2=2.5128980664188426\times10^{-3}\ \mathrm{eV}^2.
\]
These anchored absolute numbers are compare-only because the live theorem branch still lacks the
single positive normalization scalar that would emit the overall neutrino mass scale without an
external \(\Delta m_{32}^2\) anchor. Concretely, the remaining neutrino object is one
\(\lambda_\nu>0\) such that
\[
m_i=\lambda_\nu \widehat m_i,
\qquad
\Delta m_{ij}^2=\lambda_\nu^2 \widehat{\Delta}_{ij}.
\]
The current weighted-cycle branch therefore already fixes the PMNS data, the hierarchy ratio, and the
scale-free mass normal form; it misses only the single positive scalar that would place that
normal form on an absolute eV scale without an external oscillation anchor.

On the same-label neutrino-only branch, the centered eta-class is exactly \(S_3\)-isotropic, so
the same-label data are edge-constant and the solar \(1\!-\!2\) split cannot open there by itself.
The first honest solar mover therefore has to come from realized flavor-side same-label gap/defect
readback, not from another neutrino-only selector tweak.

\subsection{Intrinsic neutrino eta-chain}

The proof-facing input is smaller than the raw same-label matrix payload. It compresses to the
same-label scalar certificate
\[
\bigl(g_e,\ \omega_e\bigr)_{e\in\{12,23,31\}},
\qquad
\omega_e=\mathrm{same\text{-}label\ overlap}_e^2,
\qquad
d_e=1-\omega_e,
\]
modulo one common scale. From that certificate one forms
\[
q_e=\sqrt{g_e d_e},
\qquad
\eta_e=\log q_e-\frac13\sum_f \log q_f,
\qquad
e\in\{12,23,31\},
\]
the intrinsic selector depends only on the centered class
\[
[\eta_e]\in \mathbb{R}^3/\langle(1,1,1)\rangle.
\]
Equivalently one may work with any positive normalized family
\[
\mu_e=\frac{e^{\eta_e}}{\frac13\sum_f e^{\eta_f}}.
\]
Common scaling cancels identically, so the intrinsic selector is determined by the centered
eta-class alone.

On the isotropic neutrino-only branch one has
\[
(\eta_{12},\eta_{23},\eta_{31})=(0,0,0),
\]
which emits the intrinsic masses
\[
(m_1,m_2,m_3)=
(2.3986448447627196,\ 2.3986448447627196,\ 2.590074050773907)\times10^{-12}\ \mathrm{GeV},
\]
with
\[
\Delta m_{21}^2=0,
\qquad
\Delta m_{31}^2=9.549864971855843\times10^{-25}\ \mathrm{GeV}^2.
\]
This is the exact neutrino-only isotropy obstruction behind the remaining solar-splitting gap.

\begin{theorem}[Same-label scalar certificate suffices for intrinsic neutrino masses]
Assume the same-label gap and overlap scalars are emitted on all realized arrows. Then the full
intrinsic neutrino mass-eigenstate bundle factors through the scalar certificate
\[
\bigl(g_e,\omega_e\bigr)_{e\in\{12,23,31\}}
\]
or equivalently through the centered eta-class \([\eta_e]\). No additional raw matrix payload is
needed to form the intrinsic selector, the depressed-cubic spectrum, the ordered splittings, or
the intrinsic neutrino-side basis.
\end{theorem}

\begin{theorem}[Exact principal-branch selector from the centered eta-class]
Fix the cycle sum
\[
\Omega=\psi_{12}+\psi_{23}+\psi_{31},
\qquad |\Omega|<\frac{3\pi}{2},
\]
and a positive weight family \(\mu_e\). On the principal branch
\(\psi_e\in(-\pi/2,\pi/2)\), the affine energy
\[
A(\psi)=\sum_e \mu_e(1-\cos\psi_e)
\]
has a unique minimizer. It is given by
\[
\psi_e^\star=\arcsin(\lambda/\mu_e),
\]
where \(\lambda\) is the unique solution of
\[
\sum_e \arcsin(\lambda/\mu_e)=\Omega,
\qquad
\lambda\in(-\min_e\mu_e,\min_e\mu_e).
\]
\end{theorem}

\begin{proof}
The Euler--Lagrange equations are \(\mu_e\sin\psi_e=\lambda\). On the principal branch the scalar
function
\[
F(\lambda)=\sum_e\arcsin(\lambda/\mu_e)
\]
is strictly increasing because
\[
F'(\lambda)=\sum_e\frac{1}{\sqrt{\mu_e^2-\lambda^2}}>0.
\]
Its range is \(( -3\pi/2,3\pi/2)\), so \(F(\lambda)=\Omega\) has a unique solution. Strict
convexity follows because the Hessian of \(A\) is
\[
\nabla^2A=\mathrm{diag}(\mu_e\cos\psi_e),
\]
which is positive definite on the principal branch and remains positive definite after
restriction to the affine plane \(\sum_e\psi_e=\Omega\).
\end{proof}

Let
\[
a=m_\star=\frac{v^2}{\mu_u},
\qquad
\rho = |(M_0)_{12}|.
\]
Then the intrinsic Majorana matrix is
\[
M(\eta)=
\begin{pmatrix}
a & \rho e^{i\psi_{12}} & \rho e^{i\psi_{31}}\\
\rho e^{i\psi_{12}} & a & \rho e^{i\psi_{23}}\\
\rho e^{i\psi_{31}} & \rho e^{i\psi_{23}} & a
\end{pmatrix}.
\]

\begin{theorem}[Exact intrinsic spectral cubic]
Let \(H=M^\dagger M\). Then
\[
H=dI+T,
\qquad
d=a^2+2\rho^2,
\]
where \(T\) has zero diagonal and off-diagonals
\[
x_{12}=2a\rho\cos\psi_{12}+\rho^2e^{i(\psi_{23}-\psi_{31})},
\]
\[
x_{23}=2a\rho\cos\psi_{23}+\rho^2e^{i(\psi_{31}-\psi_{12})},
\]
\[
x_{13}=2a\rho\cos\psi_{31}+\rho^2e^{i(\psi_{12}-\psi_{23})}.
\]
Define
\[
P=|x_{12}|^2+|x_{23}|^2+|x_{13}|^2,
\qquad
Q=\Re\!\bigl(x_{12}x_{23}\overline{x_{13}}\bigr).
\]
Then the eigenvalues \(\lambda_k\) of \(T\) are exactly the three real roots of
\[
\lambda^3-P\lambda-2Q=0,
\]
and the intrinsic squared masses are
\[
m_k^2=d+\lambda_k.
\]
Equivalently,
\[
\lambda_k=
2\sqrt{\frac{P}{3}}
\cos\!\left(
\frac13\arccos\!\left(\frac{3\sqrt3\,Q}{P^{3/2}}\right)-\frac{2\pi k}{3}
\right).
\]
\end{theorem}

\begin{corollary}[Intrinsic eta-chain closure]
Once the centered eta-class is emitted at the flavor boundary, the intrinsic neutrino branch
emits
\[
\nu_1,\nu_2,\nu_3,\qquad
\Delta m_{21}^2,\Delta m_{31}^2,\Delta m_{32}^2,
\]
the ordering, and the intrinsic neutrino-side basis \(U_\nu\) with no PMNS import and no
flavor-label leakage.
\end{corollary}

\subsection{Perturbative laws around the isotropic point}

Write
\[
\psi_e=\varphi+\delta_e,
\qquad
\delta_{12}+\delta_{23}+\delta_{31}=0.
\]
At first order in the centered eta-class,
\[
\delta_e=-\tan\varphi\,\eta_e+O(\eta^2).
\]
Define
\[
\sigma^2=\frac23\sum_e \delta_e^2.
\]
Then
\[
m_{1,2}^2=m_d^2\mp 2a\rho\sin\varphi\,\sigma+O(\delta^2),
\]
so
\[
\Delta m_{21}^2 = 4a\rho\sin\varphi\,\sigma+O(\delta^2)
=4a\rho\frac{\sin^2\varphi}{\cos\varphi}
\sqrt{\frac23\sum_e\eta_e^2}+O(\eta^2).
\]

The ordered atmospheric gaps obey the corrected first-order laws
\[
\Delta m_{31}^2 = \Delta_{\mathrm{atm,iso}} + \frac12\Delta m_{21}^2 + O(\delta^2),
\qquad
\Delta m_{32}^2 = \Delta_{\mathrm{atm,iso}} - \frac12\Delta m_{21}^2 + O(\delta^2),
\]
so the first-order invariant atmospheric object is the collective-to-doublet centroid gap
\[
\Delta_{\mathrm{cent}}=
m_3^2-\frac12(m_1^2+m_2^2)
=
\Delta_{\mathrm{atm,iso}}+O(\delta^2).
\]
Its quadratic shift is
\[
\delta\Delta_{\mathrm{cent}}
=
-
a\rho\,\sigma^2
\frac{a(4\cos^2\varphi-1)+6\rho\cos\varphi}
     {2a\cos\varphi+\rho}
+O(\delta^3).
\]

The largest-mass singular vector obeys the first-order deformation law
\[
u_3=
u+
\kappa
\begin{pmatrix}
\delta_{23}\\
\delta_{31}\\
\delta_{12}
\end{pmatrix}
+O(\delta^2),
\qquad
u=\frac1{\sqrt3}(1,1,1)^T,
\qquad
\kappa=
\frac{\sqrt3\,(2a\sin\varphi+3i\rho)}
     {9(2a\cos\varphi+\rho)}.
\]
Equivalently,
\[
u_3=
u-
\tan\varphi\,\kappa
\begin{pmatrix}
\eta_{23}\\
\eta_{31}\\
\eta_{12}
\end{pmatrix}
+O(\eta^2).
\]

One exact demonstration uses centered eta-class
\[
(\eta_{12},\eta_{23},\eta_{31})
=
(0.13631072512014578,\,-0.18362987264261327,\,0.0473191475224675),
\]
and emits intrinsic masses
\[
(m_1,m_2,m_3)=
(2.3929601069646055,\,
2.4048200109774875,\,
2.589606227283229)\times10^{-12}\ \mathrm{GeV},
\]
with
\[
\Delta m_{21}^2 = 5.690121167370743\times10^{-26}\ \mathrm{GeV}^2,
\qquad
\Delta m_{31}^2 = 9.798023388600230\times10^{-25}\ \mathrm{GeV}^2.
\]
These are exact outputs of the intrinsic eta-chain once the centered eta-class is supplied. They
are not yet flavor-labeled OPH rows.

\subsection{Honest remaining blockers}

The remaining task is not ``emit PMNS from the shared basis,'' and it is not
``repair the PMNS branch.'' The weighted-cycle route already closes those steps.
The real remaining burden is to emit the one positive normalization scalar that turns the current
dimensionless hierarchy into theorem-grade absolute masses and absolute splittings. No hidden
discrete branch remains on that lane: the remaining freedom is only the positive
rescaling orbit of the absolute spectrum. Until that scalar is emitted, the weighted-cycle branch
gives a physical PMNS/hierarchy pattern and a compare-only absolute spectrum once one
atmospheric anchor is supplied.

\section{Hadrons, QCD, and the Emergence of Ordinary Matter}

The hadron derivation is the most operationally demanding part of the whole particle story. The
honest frontier is not a missing symbolic theorem but executed unquenched production data plus
production statistical and systematic control. Accordingly, the current hadron rows are not
derived outputs of this paper. The derivation nevertheless has a complete
mathematical execution bridge, a deterministic runtime receipt, a full writeback/evaluation path,
and a surrogate end-to-end executable bundle that closes the declared execution surface without
pretending to be production lattice QCD.

\subsection{Seeded \(2+1\) family and unquenched measure}

Let
\[
m_l := \frac{m_u+m_d}{2},
\qquad
\rho_l := \frac{m_u+m_d}{2\,\Lambda_{\overline{\mathrm{MS}}}^{(3)}},
\qquad
\rho_s := \frac{m_s}{\Lambda_{\overline{\mathrm{MS}}}^{(3)}}.
\]
The seeded fixed-physics family is
\[
a\Lambda_n = a\Lambda_{\mathrm{seed}}\,2^{-n},
\qquad
a\Lambda_{\mathrm{seed}} = \sqrt{\rho_l \rho_s},
\]
\[
a m_l^{(n)} = a\Lambda_n \rho_l,
\qquad
a m_s^{(n)} = a\Lambda_n \rho_s,
\]
\[
\beta_n = 6 + \frac{9}{2\pi^2}\log\!\frac{1}{a\Lambda_n},
\qquad
L_n = \left\lceil \frac{\lambda_L^{\mathrm{target}}}{a\Lambda_n}\right\rceil,
\qquad
T_n = \left\lceil \frac{\lambda_T^{\mathrm{target}}}{a\Lambda_n}\right\rceil.
\]
The unquenched ensemble measure is
\[
d\mu_n(U)
=
Z_n^{-1}
\exp[-S_g(U;\beta_n)]
\det D_l(U; a m_l^{(n)})^2
\det D_s(U; a m_s^{(n)})
\,dU.
\]
On the present branch, \(N_f=2+1\) and QED is off.

\subsection{Runtime receipt and deterministic cfg/source contract}

The emitted execution law is
\[
U_{n,c}
=
K_n^{N_{\mathrm{therm}} + (\mathrm{cfg\_index})\,N_{\mathrm{sep}}}(U_{\mathrm{cold}}; \mathrm{seed}_{n,c}),
\]
with stop-time formula
\[
t_{\mathrm{stop}}(n,c)
=
N_{\mathrm{therm}} + (\mathrm{cfg\_index})\,N_{\mathrm{sep}}.
\]
The deterministic cfg seed law is
\[
\mathrm{seed}_{n,c}
=
\mathrm{bytes.fromhex}(\mathrm{cfg\_seed\_hash}_{n,c}),
\]
with
\[
\mathrm{cfg\_seed\_hash}
=
\mathrm{SHA256}\!\Big(
\mathrm{Serialize}\!\big(
\mathrm{ensemble\_id},\beta,L,T,a m_l,a m_s,\mathrm{cfg\_index}
\big)
\Big).
\]
The source set is fixed as
\[
S_{n,c} = \{ [0,0,0,0],\ [L_n//2, L_n//2, L_n//2, T_n//2] \}.
\]
The non-null external runtime receipt used in the surrogate execution is
\[
N_{\mathrm{therm}} = 2048,
\qquad
N_{\mathrm{sep}} = 512.
\]
These values are surrogate execution inputs only; they are not claimed as theorem outputs.

The production geometry summary makes the remaining runtime boundary concrete. On the emitted
seeded \(2+1\) family there are three ensembles and six configurations total. If one stores all
four links at every site as full double-complex \(3\times 3\) matrices, the naive raw gauge
storage estimate is
\[
576\ \mathrm{bytes/site},
\]
which gives total naive raw gauge storage
\[
2.80071464105088\times 10^{14}\ \mathrm{bytes}
\]
over the present production schedule. By contrast, the backend correlator dump needed by the
analysis layer is tiny:
\[
195264\ \mathrm{bytes}
\]
for the full \(\pi_{\mathrm{iso}}\), \(N_{\mathrm{iso,dir}}\), and \(N_{\mathrm{iso,ex}}\)
payload on the present schedule. So the live hadron residual is not the analysis schema but the
actual backend correlator dump from real production execution.

\subsection{Operational reading of the remaining computation}

It is useful to separate the lightweight analysis layer from the expensive numerical layer.
Informally, the repository-side control plane is already present: the seeded ensemble family, the
deterministic cfg/source contract, the export manifest, the jackknife evaluator, the forward-window
selector, and the publication-budget schema are all explicit. The missing step is therefore not
another bookkeeping script but the physical generation of dynamical gauge configurations and
measured correlator arrays on the fixed \(N_f=2+1\), QED-off branch.

Technically, the remaining production computation is the standard lattice-QCD stable-channel
workflow on the emitted family: run unquenched RHMC/HMC for the three seeded ensembles; for each
realized cfg and fixed source solve the clover-improved Wilson light/strange systems; construct the
zero-momentum \(\pi_{\mathrm{iso}}\), \(N_{\mathrm{iso,dir}}\), and \(N_{\mathrm{iso,ex}}\)
two-point sequences; write the backend bundle; convert that bundle into the repo-side production
dump; then let the existing jackknife and forward-window machinery emit
\(a m_{X,\mathrm{ground}}\), \(R_X\), \(m_X[\mathrm{GeV}]\), and the published
\(\sigma_{\mathrm{stat}}\), \(\delta_{\mathrm{cont}}\), \(\delta_{\mathrm{vol}}\), and
\(\delta_{\chi}\) fields. In that sense the hadron lane is analysis-complete but
execution-incomplete.

The engineering burden is therefore asymmetric. Local execution is sufficient for manifest
generation, surrogate validation, and downstream readout tests, because those stages operate on a
tiny correlator payload. The physical branch is different: the dominant cost is upstream
gauge-field generation and Dirac solves, not the final analysis dump. This is therefore a classical
high-performance-computing workload, not a special quantum-computing shortcut. Current
quantum-cloud hardware can support toy lattice-gauge demonstrations, but it is not presently a
realistic execution target for full \(3+1\)-dimensional dynamical \(SU(3)\) lattice-QCD hadron
spectroscopy. So the honest deployment target is a multi-node HPC cluster or cloud HPC allocation
rather than a laptop, ordinary workstation, or current quantum-cloud device. As an operational estimate, a backend adapter and first real stable-channel pilot dump
are a week-scale software task, whereas publishable \(\pi_{\mathrm{iso}}/N_{\mathrm{iso}}\)
systematics are month-scale numerical work and depend strongly on the available allocation, solver
quality, and whether one stays on the current stable-channel branch or also opens the later
\(\rho\)-scattering and isospin-breaking program.

\subsection{Stable-channel correlators, effective masses, and forward windows}

The pion stable channel is
\[
p_\pi^{(n,c,s)}(t)
=
\sum_x \Re\,\mathrm{tr}_{c,\mathrm{spin}}
\left[
\gamma_5 S_l(x;s)\gamma_5 S_l(s;x)
\right].
\]
The nucleon stable channel is
\[
p_{N,\mathrm{dir}}^{(n,c,s)}(t) = \sum_x G_d,
\qquad
p_{N,\mathrm{ex}}^{(n,c,s)}(t) = \sum_x G_x,
\]
\[
p_N^{(n,c,s)}(t)
=
p_{N,\mathrm{dir}}^{(n,c,s)}(t) - p_{N,\mathrm{ex}}^{(n,c,s)}(t).
\]
Cfg/source and ensemble averaging are
\[
\bar p_X^{(n,c)}(t)
=
\frac{1}{|S_{n,c}|}\sum_{s\in S_{n,c}} p_X^{(n,c,s)}(t),
\]
\[
C_X^{(n)}(t)
=
\frac{1}{|C_n|}\sum_{c\in C_n} \bar p_X^{(n,c)}(t).
\]

The effective-mass laws are
\[
a m_{\mathrm{eff},\pi}(t)
=
\log\!\frac{C_\pi(t)}{C_\pi(t+1)},
\qquad
a m_{\mathrm{eff},N}(t)
=
\log\!\frac{|C_N(t)|}{|C_N(t+1)|}.
\]
The forward window is
\[
W_n = \{ t : 1 \le t+1 < \lfloor T_n/2 \rfloor \}.
\]
The evaluator monitors log-convexity,
\[
R_{\log\mathrm{conv},\pi}(t)
=
C_\pi(t)^2 - C_\pi(t-1)C_\pi(t+1),
\]
\[
R_{\log\mathrm{conv},N}(t)
=
|C_N(t)|^2 - |C_N(t-1)|\,|C_N(t+1)|,
\]
tail-drop,
\[
D_X(t)
=
a m_{\mathrm{eff},X}(t) - a m_{\mathrm{eff},X}(t+1),
\]
and mirror suppression,
\[
M_X(t)
=
\exp[-a m_{\mathrm{eff},X}(t)(T_n - 2t)].
\]
The selected forward window is the longest contiguous run satisfying finite effective masses,
nonnegative log-convexity up to tolerance, nonnegative tail-drop up to tolerance, mirror
suppression below threshold, and local plateau flatness. The candidate ground-state mass is then
the weighted window average
\[
a m_{X,\mathrm{ground}}
=
\frac{\sum_{t\in W_X^{\mathrm{sel}}} w_t\,a m_{\mathrm{eff},X}(t)}
{\sum_{t\in W_X^{\mathrm{sel}}} w_t},
\qquad
w_t = \frac{1}{\max(\sigma_t^2,\varepsilon)}.
\]

\subsection{Statistics, systematics, and dimensional readout}

Delete-1 jackknife is performed over the cfg axis after source averaging inside each cfg. With
\(n_{\mathrm{cfg}}\) configurations and integrated autocorrelation time
\(\tau_{\mathrm{int},\mathrm{cfg}}\), the effective cfg count is
\[
n_{\mathrm{eff},\mathrm{cfg}}
=
\frac{n_{\mathrm{cfg}}}{2\,\tau_{\mathrm{int},\mathrm{cfg}}}.
\]
The published statistical error is
\[
\sigma_{\mathrm{stat},X}
=
\mathrm{JKstderr}(a m_{X,\mathrm{ground}}).
\]
The machine-readable systematics field uses
\[
\sigma_{\mathrm{sys},X}
=
\sqrt{
\delta_{\mathrm{cont},X}^2
+
\delta_{\mathrm{vol},X}^2
+
\delta_{\chi,X}^2
}.
\]
Continuum, volume, and chiral proxies are encoded as
\[
R_X^{(n)} = \frac{a m_X^{(n)}}{a\Lambda_n},
\qquad
R_X^{(n)} \approx R_X(0) + c_X (a\Lambda_n)^2,
\]
\[
\delta_{\mathrm{cont},X}^{(n)}
=
a\Lambda_n\,\left|R_X^{(n)} - R_X(0)\right|,
\]
\[
\delta_{\mathrm{vol},\pi}^{(n)}
=
a m_\pi^{(n)}
\frac{e^{-a m_\pi^{(n)}L_n}}{\max(a m_\pi^{(n)}L_n,1)},
\]
\[
\delta_{\mathrm{vol},N}^{(n)}
=
a m_N^{(n)}
\frac{e^{-a m_\pi^{(n)}L_n}}{\max(a m_\pi^{(n)}L_n,1)},
\]
\[
\delta_{\chi,\pi}^{(n)}
=
a\Lambda_n \left|R_\pi^{(n)} - \langle R_\pi\rangle_n\right|,
\]
\[
Q_N^{(n)} = \frac{R_N^{(n)}}{R_\pi^{(n)}},
\qquad
\delta_{\chi,N}^{(n)}
=
a\Lambda_n \langle R_\pi\rangle_n
\left|Q_N^{(n)} - \langle Q_N\rangle_n\right|.
\]
These are surrogate publication proxies, not production physical systematics. Given a
ground-state candidate,
\[
R_X = \frac{a m_{X,\mathrm{ground}}}{a\Lambda_{\overline{\mathrm{MS}}}^{(3)}},
\qquad
m_X[\mathrm{GeV}] = R_X\,\Lambda_{\overline{\mathrm{MS}}}^{(3)}[\mathrm{GeV}],
\]
with
\[
\Lambda_{\overline{\mathrm{MS}}}^{(3)} = 0.3344017073\ \mathrm{GeV}.
\]

\subsection{Surrogate execution bundle and honest frontier}

The surrogate execution evolves a latent state \(z\in\mathbb{R}^d\) with Hamiltonian
\[
H(z,p) = S(z) + \frac12 \sum_i p_i^2,
\]
where
\[
S(z)
=
\frac12 \sum_i \omega_i^2 z_i^2
+
\lambda_4 \sum_i z_i^4
+
\kappa \sum_i (z_{i+1}-z_i)^2
+
2\alpha_l \sum_i \log(\mu_l^2 + z_i^2)
+
\alpha_s \sum_i \log(\mu_s^2 + \tfrac12 z_i^2).
\]
Leapfrog integration plus Metropolis accept/reject gives the surrogate HMC update
\[
(z,p) \mapsto (z',p'),
\qquad
P_{\mathrm{acc}} = \min(1,e^{-\Delta H}).
\]
This kernel is not physical lattice-QCD RHMC/HMC. It is a deterministic executable surrogate
honoring the emitted receipt and seed law.

For validation of the execution bridge, the surrogate locks the ground-state masses to the hadron
audit proxy values
\[
m_{\pi,\mathrm{proxy}} = 0.13497682776768472\ \mathrm{GeV},
\qquad
m_{N,\mathrm{iso,proxy}}
=
\frac{m_p+m_n}{2}
=
0.93891875434\ \mathrm{GeV}.
\]
Thus
\[
a m_{\pi,\mathrm{proxy}}^{(n)}
=
\frac{m_{\pi,\mathrm{proxy}}}{\Lambda_{\overline{\mathrm{MS}}}^{(3)}[\mathrm{GeV}]}
\,a\Lambda_n,
\]
\[
a m_{N,\mathrm{iso,proxy}}^{(n)}
=
\frac{m_{N,\mathrm{iso,proxy}}}{\Lambda_{\overline{\mathrm{MS}}}^{(3)}[\mathrm{GeV}]}
\,a\Lambda_n.
\]
The surrogate correlator is
\[
C_X^{\mathrm{sur}}(t)
=
A_{X,0} e^{-a m_X t}
+
A_{X,1} e^{-a m_{X,\mathrm{ex}} t}
+
A_{X,\mathrm{mir}} e^{-a m_X (T_n-t)}
\]
up to a multiplicative correlated noise factor. Direct and exchange nucleon pieces are written as
\[
C_{N,\mathrm{dir}}^{\mathrm{sur}}(t)
=
f_{\mathrm{dir}}\,C_N^{\mathrm{sur}}(t),
\qquad
C_{N,\mathrm{ex}}^{\mathrm{sur}}(t)
=
(f_{\mathrm{dir}}-1)\,C_N^{\mathrm{sur}}(t),
\]
so that
\[
C_N^{\mathrm{sur}}(t)
=
C_{N,\mathrm{dir}}^{\mathrm{sur}}(t)
-
C_{N,\mathrm{ex}}^{\mathrm{sur}}(t).
\]

On the finest surrogate ensemble, the resulting diagnostic candidates are
\[
m_{\pi,\mathrm{iso}}^{\mathrm{sur}} = 0.135039383836\ \mathrm{GeV},
\qquad
m_{N,\mathrm{iso}}^{\mathrm{sur}} = 0.938960210578\ \mathrm{GeV},
\]
with worst absolute error approximately \(6.26\times10^{-5}\ \mathrm{GeV}\) across the surrogate
stable-channel family. These numbers validate the emitted execution bridge. They are not
promotable production hadron predictions.

So the hadron derivation closes the full
\[
\text{receipt} \to \text{execution} \to \text{writeback} \to \text{evaluation} \to \text{budgets} \to \text{forward-window summary}
\]
path on executed surrogate data, while the remaining physical closure requires production
unquenched RHMC/HMC, real Dirac solves and baryon contractions, production autocorrelation
studies, and production continuum / finite-volume / chiral systematics.

\section{Observer-Centric Particle Ontology and Measurement}

The particle paper cannot stop at masses and couplings. OPH is observer-centric at the level of
its basic ontology, so a particle chapter also has to explain what a particle \emph{is} in this
language and how a measurement turns an excitation surface into an actual observed state.
Fortunately, this is one of the places where the OPH theorem surface contains
real mathematical content rather than only interpretation.

\subsection{Observer patches, overlap consistency, and particle data}

The consensus paper formulates the basic kinematic picture in patch-net language: each observer
patch carries a local state, neighboring patches compare a shared interface alphabet, and a
global state is physically admissible exactly when neighboring projections agree on the overlap.
The algebraic OPH version of the same statement is Axioms~\ref{ax:screen} and
\ref{ax:overlap}: local physical data are carried by patch algebras \(\mathcal A(P)\), and those local
states must agree on shared subalgebras whenever patches overlap.

This matters for particles because OPH does not begin with a single absolute global particle
basis. It begins with patch-local algebras, overlap-visible observables, edge sectors, and
transport data. A particle row in the reported derivation is therefore not best
understood as an isolated primitive label. It is a readout claim about a stable or
candidate-stable excitation structure visible to a family of observer patches and consistent on
their overlaps.

In the observer formulation, an observer is represented by
\[
O=(P,\mathcal A(P),\rho,R),
\]
where \(P\) is the observer patch, \(\mathcal A(P)\) its local algebra, \(\rho\) the local
state, and \(R\) the record algebra. The key point for the particle derivation is that \(R\) is not
an extra classical layer stapled onto quantum theory. It is built from approximately commuting
projectors in overlap-visible centers and record registers, so it is shareable across overlapping
observer descriptions without violating the usual no-cloning constraints on generic quantum
states.

\subsection{Record algebras and definite outcomes}

The technical supplement and the screen-microphysics paper make the measurement
interface precise at fixed cutoff. A record in patch \(P\) is a family of orthogonal projectors
\(\{\Pi_r\}\subset\mathcal A(P)\) whose generated algebra is approximately abelian, dynamically
stable, and shareable through overlap centers. On the explicit fixed-cutoff screen architecture,
the record algebra after a completed write/verify cycle is
\[
\mathcal Z_{\mathrm{rec}}(t)
=
\mathrm{Alg}\Bigl(
\{P_{m_I^{(J)}(t)}\}_{I\neq J},
\{P_{r_I^{\mathrm{bulk}}(t)}\}_I,
\{\Pi_\alpha^{(IJ)}(t)\}_{I<J,\alpha}
\Bigr).
\]
The microphysics paper proves that, on the declared operational measurement surface,
\(\mathcal Z_{\mathrm{rec}}(t)\) is commutative and central for the readout instrument being
used.

This fixed-cutoff centrality result is what lets OPH talk about definite outcomes without adding
a second ontology. The technical supplement phrases the same idea through edge-center
decomposition: on a fixed collar with exact Markov structure, the state splits into
superselection blocks,
\[
\rho_{ABC}
=
\bigoplus_j
q_j\,
\rho^{(j)}_{A b_L^{(j)}}\otimes \rho^{(j)}_{b_R^{(j)} C},
\]
and the label \(j\) is classical center data. The supplement then goes one step further: in the
physical algebra there are no interference observables between different sector blocks. So once a
particle detection event has been recorded in the observer-accessible center/record algebra, the
accessible physics is organized as a classical mixture over those blocks rather than as a
superposed measurement record.

For the particle paper, this means a measured particle family or excitation channel should be
thought of as an observer-accessible sector or record value carried by the central measurement
surface, not as a mysterious absolute collapse of the full universe-state into a preferred basis
chosen by hand.

\subsection{Born probabilities and post-measurement states}

The fixed-cutoff measurement package is also explicit about probabilities and updates. For every
event \(E\) in the sigma algebra generated by \(\mathcal Z_{\mathrm{rec}}(t)\), the measurement
probability is
\[
\mathbb P_t(E)=\operatorname{Tr}\!\bigl(\rho_t P_E\bigr),
\]
where \(P_E\in\mathcal Z_{\mathrm{rec}}(t)\) is the projector for that event. Conditioning on the
event then produces the post-measurement state
\[
\rho_t\!\mid_E
=
\frac{P_E\rho_t P_E}{\operatorname{Tr}(\rho_t P_E)}.
\]
In the supplement, the same statement appears as the L\"uders update on the record algebra. In
the observer wrapper, this fixed-cutoff Born/L\"uders package is imported as theorem-bearing
rather than interpretive rhetoric.

This directly answers the particle-state question in OPH language. Measurements give particles
actual states by conditioning the observer-accessible state on the central record algebra. A
detector click, a stable overlap-sector readout, or a pointer value does not merely announce a
pre-existing classical label; it defines the conditioned state for subsequent observer-accessible
physics. If the same event is re-read without any accepted repair move touching its support, the
microphysics paper proves that the re-read probability is \(1\). So the measurement interface is
not only well defined but operationally auditable.

\subsection{Particle identity as transport-stable structure}

Once one asks not only ``which event happened?'' but also ``which particle family is this?'', the
flavor derivation becomes essential. The active flavor derivation does not start with the
names electron, muon, tau lepton, or up quark built into the fundamental observable. It starts
with transport kernels, generation-bundle data, same-label eigenline transport, overlap-edge
cocycles, and a persistent flavor observable carrying intrinsic labels \(f1,f2,f3\). This is an
important honesty condition in the paper. At the deepest currently active flavor surface, family
identity is transport-stable intrinsic structure first and named experimental family assignment
second.

That is why the Yukawa/excitation dictionary remains such a central open task. The flavor
pipeline constructs a substantial invariant object: it exports projectors, spectral
gaps, pair suppressions, cycle phases, and common sector-response objects. But the map from those
intrinsic structures to the named low-energy fermion families is exactly part of the unfinished
closure problem. The quark derivation emits named rows because the forward
branch is far enough along to support direct numerical comparison there; the charged-lepton and
neutrino derivations do not yet do so quantitatively, even though they use the same downstream family
machinery.

So in the strongest justified sense, particle identity in OPH is a refinement-stable
observer-accessible transport pattern whose named interpretation is inherited from the
available dictionary and comparison surfaces. Where that dictionary is still open, the paper should say
so explicitly.

\subsection{Scope boundary of the measurement chapter}

This whole chapter has to preserve one final distinction. The fixed-cutoff measurement interface
is theorem-bearing on the microphysics surface: central record algebra, Born
probabilities, L\"uders conditioning, and repeated-read stability are written there as
fixed-cutoff statements. What is \emph{not} yet proved at the same level is the stronger global
observer-continuation or strange-loop story sometimes discussed in broader OPH wrapper notes. The
particle paper should therefore use the fixed-cutoff measurement package directly, while leaving
the larger metaphysical closure story out of the main technical chain.

% Later sections to be filled after source extraction and status-table synchronization.
% \section{Results at a Glance}
\section{Discussion: Matter, Antimatter, Supersymmetry, and Other Open Questions}

This section is intentionally not written in the same voice as the structural and quantitative
sections above. The OPH theorem surface distinguishes sharply between recovered-core
results, calibration branches, secondary quantitative branches, and continuations. Matter versus
antimatter, supersymmetry, and several adjacent questions lie on that boundary. They are too
important to omit from a particle-spectrum paper, but they are not yet part of the finished
prediction surface in the same way that the bosonic and quark rows are.

\subsection{Matter and antimatter}

The Standard Model structural branch fixes the chiral gauge architecture within which
matter and antimatter are defined. Once a chiral fermion representation is present, its conjugate
representation gives the corresponding antiparticle content. In that limited sense, OPH explains
why matter and antimatter both exist: they are paired by the realized gauge and chiral structure
of the low-energy branch.

What is \emph{not} yet structurally closed is the cosmological asymmetry between them. The
technical supplement records a continuation-level baryogenesis story tied to the realized
\(\mathbb Z_6\) quotient and the electroweak topological channel. In that continuation, the
entropy deficit of the \(\mathbb Z_6\) quotient gives a suppression factor
\[
\varepsilon = e^{-\ln 6} = \frac{1}{6},
\]
and combining this with the \(12\)-fermion electroweak baryon-violating channel gives the
order-of-magnitude estimate
\[
\varepsilon^{12}=6^{-12}\approx 4.6\times 10^{-10},
\]
which is numerically close to the observed baryon asymmetry scale.

But the observer-side status notes are explicit about the status of that estimate: the baryon
asymmetry scale remains a Phase-III deferred continuation, and suppression-counting estimates are
not a substitute for a derived out-of-equilibrium mechanism. The present derivation
also does not expose a baryogenesis output row. So the particle paper can report the
existing OPH continuation idea, but it should not present a sharp matter/antimatter abundance
prediction as if it were a closed theorem or an emitted output.

\subsection{Why the present branch does not require supersymmetry}

The technical supplement contains a dedicated section on supersymmetry and gauge-coupling
unification. Its main point is not that supersymmetry has been mathematically ruled out in every
possible OPH extension. Its point is narrower and more relevant to this paper: the present
OPH branch does not \emph{need} a supersymmetric partner spectrum in order to explain why
unification-like running might appear.

At one loop, the supplement notes the familiar fact that Standard Model beta-function
coefficients do not produce successful naive unification, whereas MSSM-like coefficients do. The
OPH continuation idea is then that edge-sector heat-kernel weights can shift the effective
beta-function coefficients in an MSSM-like direction through geometric or entropic sector
multiplicity rather than through an actual low-energy superpartner spectrum. In that reading,
``unification-like behavior'' and ``a supersymmetric particle zoo'' come apart.

For the particle paper, the correct conclusion is therefore conditional and modest:
\begin{enumerate}[leftmargin=2em]
\item the OPH derivation reconstructs the realized Standard Model branch without having to
introduce a supersymmetric partner for every known field;
\item the supplement contains a continuation-level argument that some unification-style running
features could arise from edge-sector structure rather than MSSM particle content;
\item none of this is yet a universal theorem that supersymmetry is impossible.
\end{enumerate}

So if the paper asks ``why is there no supersymmetry in the derived spectrum?'', the best
technical answer is: because the realized OPH branch stops at the Standard
Model content, and the existing unification discussion is trying to reproduce the relevant
running behavior geometrically rather than by extending the particle content to a full
supersymmetric multiplet structure.

\subsection{Three generations, flavor closure, and what remains open}

The structural Standard Model branch fixes \(N_g=3\) on the realized MAR-admissible
branch. But the particle paper should not blur this result together with the much harder flavor
closure problem. Three generations are structurally recovered; the full flavor dictionary,
excitation map, CKM closure, charged-lepton closure, and neutrino closure remain partially open.
This is exactly why the flavor chapter and the later family chapters spend so much time on
constructive object boundaries rather than only on final numbers.

\subsection{Ordinary matter and the missing hadron closure}

A second distinction that belongs in discussion form rather than in the results sections is the
difference between ``elementary rows exist'' and ``ordinary matter is fully derived.'' Ordinary
visible matter is dominated by hadrons, especially protons and neutrons. The present OPH
derivation has a serious hadron pipeline, but it still lacks promotable hadron rows on the
default surface. So the paper can honestly claim that the route toward ordinary matter is
architecturally present, but it cannot yet claim that the full observed matter spectrum has been
numerically derived end to end.

\subsection{How to read these open questions}

The reader should therefore interpret this whole discussion section with the same rule used by
\emph{Recovering Relativity and Standard Model Structure from Observer Overlap Consistency}~\cite{ophcompact}. A failure in one of these continuation stories would not by itself undo the
recovered-core gauge and gravity chain, nor would it erase the reported bosonic and quark
rows. Conversely, a future success in these continuation chapters would materially strengthen the
particle-spectrum derivation without changing the logical fact that they sit beyond the closed
structural core.

\section{Conclusion}

The shortest honest summary of this draft is this. OPH has a nontrivial
particle-spectrum derivation that is much more concrete than a qualitative aspiration, but it is not
yet a uniformly closed derivation of the full particle zoo. The structural theorem surface fixes
the realized Standard Model branch and the exact massless carrier sectors. The derivation
then emits quantitative bosonic rows through electroweak calibration and the Higgs/top critical
stage, and it emits continuation-level quark rows on the flavor branch. Charged leptons,
flavor-labeled neutrinos, and hadrons each have substantial derivation chains, but those chains still
terminate at named missing constructive objects rather than at promotable masses.

That asymmetry is not something to hide; it is the research state described by the mathematics in
this draft. The paper puts the finished structural results, the emitted quantitative rows, and the
remaining closure fronts onto one common page. The visible gates remain the same ones that
organize the derivation: the Yukawa/excitation dictionary, theorem-level charged-lepton closure,
quark texture/CKM/CP closure, theorem-level neutrino closure, and the nonperturbative hadron
branch. Until those move, this paper should be read as a technical derivation document with real
outputs in hand, not as a final claim that the entire observed spectrum is finished.

\begin{thebibliography}{9}

\bibitem{ophcompact}
B.\ M\"uller,
\emph{Recovering Relativity and Standard Model Structure from Observer Overlap Consistency}.
GitHub PDF:
\url{https://github.com/FloatingPragma/observer-patch-holography/blob/main/paper/recovering_relativity_and_standard_model_structure_from_observer_overlap_consistency_compact.pdf}

\end{thebibliography}

\end{document}
